\chapter{Deep feedforward networks and backpropagation}
\label{sec:deep_feedforward_and_backpropagation}

The quest to construct artificial intelligence originated with a profound biological ambition: to emulate the computational architecture of the human brain. While classical machine learning algorithms—such as Ordinary Least Squares, Ridge regression, and linear support vector machines—operate by mapping inputs directly to predictions via fixed, handcrafted feature representations, natural sensory perception defies static linear modeling. Human cognition effortlessly navigates complex perceptual manifolds: recognizing a handwritten digit, segmenting a tumor in an acoustic scan, or parsing visual scenes across dramatic shifts in illumination, scale, and orientation.

To transcend the rigid confines of linear predictors, connectionist theory introduced the \textbf{artificial neural network}—a massively parallel computational architecture composed of elementary interconnected processing nodes. This intellectual journey began with Frank Rosenblatt's 1958 Perceptron, an elegant bio-inspired unit that computed affine combinations thresholded by a step function. Yet, as Marvin Minsky and Seymour Papert demonstrated in their devastating 1969 critique, single-layer perceptrons possess a fatal geometric blind spot: bounded strictly by linear separability, they fail on elementary non-linear predicates as simple as the Exclusive-OR (XOR) function, precipitating the First AI Winter.

The modern deep learning revolution was born from overcoming this linear barrier through two foundational mathematical innovations. First, by cascading multiple layers of neurons equipped with continuous, non-linear activation functions into \textbf{Multilayer Perceptrons (MLPs)}, networks acquire the capacity for \textbf{representation learning}. Rather than relying on human engineers to handcraft feature spaces, deep architectures automatically fold, stretch, and warp input coordinate space, transforming complex non-linear manifolds into representations that are cleanly linearly separable. The celebrated Hornik-Cybenko Universal Approximation Theorem formalizes this expressive power, guaranteeing that feedforward networks can approximate any continuous function on a compact domain to arbitrary precision.

Yet, representational power is useless without an efficient training engine. Because deep compositions of non-linear activations entangle parameters in transcendental algebraic systems, closed-form matrix inversions are mathematically impossible, and numerical finite differences require an intractable $\mathcal{O}(W^2)$ computational complexity. This computational gridlock was broken by formulating deep networks as Directed Acyclic Graphs (DAGs) and applying reverse-mode automatic differentiation: the \textbf{Backpropagation algorithm}. By propagating adjoint error sensitivities backward through the network in a single sweep, backpropagation evaluates the exact gradient across millions of parameters in linear $\mathcal{O}(W)$ time, establishing the foundational mathematical engine of contemporary artificial intelligence.

\section{The perceptron}
\label{sec:perceptron_and_linear_separability}

The quest to construct artificial intelligence originated with a profound biological ambition: to understand and replicate the computational mechanisms of the brain. In the 1940s and 1950s, the emergence of early digital computing was defined by von Neumann architectures designed for serial instruction execution and exact symbolic arithmetic. While these mechanical and electronic calculators could perform millions of numerical operations per second, they proved utterly incapable of solving elementary perceptual tasks that animal brains resolve effortlessly and instantaneously---such as recognizing a visual pattern, distinguishing spoken phonemes, or adapting to sensory variations in an uncertain environment.

This striking divergence highlighted a fundamental insight: biological intelligence does not operate via centralized instruction registers, brittle algorithmic decision trees, or hand-crafted symbolic lookup tables. Instead, cognitive function emerges from the collective dynamics of billions of interconnected, elementary biological cells---neurons---transmitting electrical impulses across plastic synaptic junctions. Despite individual biological neurons operating at millisecond timescales (orders of magnitude slower than electronic switching circuits), their massive parallel connectivity enables robust associative memory, real-time sensory processing, and the ability to learn directly from experience.

Inspired by these neurophysiological observations, early cyberneticians and cognitive scientists---beginning with Warren McCulloch and Walter Pitts in 1943, Donald Hebb in 1949, and culminating in Frank Rosenblatt in 1958---sought to abstract the physical machinery of the biological neuron into an elementary mathematical learning unit. By modeling synaptic plasticity as adjustable numerical weights and dendritic summation as an affine score thresholded by an excitation activation, they aimed to create an artificial computing system capable of discovering patterns and categorizing sensory observations autonomously without explicit human programming.

\subsection{Architecture and convergence}
\label{sec:biological_and_artificial_neuron}
\label{sec:perceptron_learning_algorithm}

\paragraph{Biological morphology and Rosenblatt's artificial threshold unit.}
In biological nervous systems, an individual neuron consists of four primary morphological structures:
\begin{itemize}[noitemsep,topsep=2pt,parsep=0pt,partopsep=0pt]
  \item \textbf{Dendrites}: Branching nerve fibers that receive incoming chemical neurotransmitter signals from adjacent upstream neurons across synaptic clefts.
  \item \textbf{Soma (Cell Body)}: The metabolic hub containing the nucleus, which integrates and accumulates incoming post-synaptic electrical potentials over space and time.
  \item \textbf{Axon}: A long, slender nerve transmission cable (often insulated by fatty myelin sheaths) that propagates an all-or-nothing action potential (voltage spike) down its length once the soma's membrane potential surpasses an excitation threshold.
  \item \textbf{Synapses}: Chemical junctions connecting axon terminal boutons to recipient dendrites, where biological synaptic plasticity modulates the efficacy of electrical transmission.
\end{itemize}

In 1958, Frank Rosenblatt~\cite{rosenblatt1958perceptron} formalized this biological mechanism into the \textbf{Perceptron}. Given an input vector of sensory features $\mathbf{x} = [x_1, \dots, x_d]^\top \in \R^d$, the artificial neuron associates each input channel with a tunable real-valued weight $w_j \in \R$, incorporates a scalar bias threshold $b \in \R$, and computes a linear pre-activation score:
\begin{equation}
z \triangleq \mathbf{w}^\top \mathbf{x} + b = \sum_{j=1}^d w_j x_j + b.
\label{eq:perceptron_preactivation}
\end{equation}
The continuous scalar score $z$ is then passed through a discontinuous Heaviside step activation function to produce a binary classification decision:
\begin{equation}
\hat{y} = f(z) = \operatorname{sgn}(z) = \begin{cases} +1 & \text{if } \mathbf{w}^\top \mathbf{x} + b \ge 0, \\ -1 & \text{if } \mathbf{w}^\top \mathbf{x} + b < 0. \end{cases}
\label{eq:perceptron_activation}
\end{equation}

\begin{figure}[H]
\centering
\begin{tikzpicture}[scale=0.92, every node/.style={transform shape}]
  % Left Panel: Biological Neuron
  \begin{scope}[shift={(0,0)}]
    \node[font=\bfseries\sffamily\color{politoNavy}] at (3.0, 2.7) {(a) Biological Neuron Anatomy};
    
    % Dendrites branching
    \draw[thick, deepblue] (0.2, 1.8) -- (1.5, 1.2);
    \draw[thick, deepblue] (0.0, 1.0) -- (1.5, 0.8);
    \draw[thick, deepblue] (0.3, 0.0) -- (1.5, 0.4);
    \draw[thick, deepblue] (-0.2, 1.5) -- (0.8, 1.2);
    \draw[thick, deepblue] (-0.2, 0.4) -- (0.8, 0.6);
    \node[font=\scriptsize\color{deepblue}, align=center] at (0.3, -0.4) {Dendrites\\(Inputs $\mathbf{x}$)};

    % Soma (cell body)
    \draw[fill=deepblue!15, draw=deepblue, very thick] (2.2, 0.8) circle (0.85);
    \draw[fill=deepblue!40, draw=deepblue] (2.2, 0.8) circle (0.35);
    \node[font=\scriptsize\bfseries\color{politoNavy}, align=center] at (2.2, 0.8) {Soma\\($\sum$)};
    \node[font=\scriptsize\color{deepblue}, align=center] at (2.2, -0.4) {Cell Body\\(Accumulator)};

    % Axon with myelin sheath
    \draw[ultra thick, deepblue] (3.05, 0.8) -- (5.2, 0.8);
    \foreach \x in {3.3, 3.9, 4.5} {
      \draw[fill=skyBlue!40, draw=politoNavy, rounded corners=2pt] (\x, 0.60) rectangle (\x+0.45, 1.00);
    }
    \node[font=\scriptsize\color{deepblue}, align=center] at (4.1, -0.4) {Axon Cable\\(Transmission)};

    % Axon Terminals / Synapses
    \draw[thick, deepblue] (5.2, 0.8) -- (5.9, 1.3);
    \draw[thick, deepblue] (5.2, 0.8) -- (6.0, 0.8);
    \draw[thick, deepblue] (5.2, 0.8) -- (5.9, 0.3);
    \foreach \pos in {(5.9, 1.3), (6.0, 0.8), (5.9, 0.3)} {
      \fill[examred] \pos circle (2.2pt);
    }
    \node[font=\scriptsize\color{examred}, align=center] at (6.0, -0.4) {Synaptic Boutons\\(Output $\hat{y}$)};
  \end{scope}

  % Vertical Separator
  \draw[dashed, gray!40, thick] (6.9, -0.8) -- (6.9, 3.1);

  % Right Panel: Artificial Perceptron
  \begin{scope}[shift={(7.8,0)}]
    \node[font=\bfseries\sffamily\color{politoNavy}] at (3.5, 2.7) {(b) Rosenblatt's Artificial Perceptron};

    % Inputs
    \node[circle, draw=deepblue, fill=deepblue!10, thick, inner sep=2pt, font=\footnotesize] (x1) at (0, 1.8) {$x_1$};
    \node[circle, draw=deepblue, fill=deepblue!10, thick, inner sep=2pt, font=\footnotesize] (x2) at (0, 0.9) {$x_2$};
    \node[font=\footnotesize] (xdots) at (0, 0.2) {$\vdots$};
    \node[circle, draw=deepblue, fill=deepblue!10, thick, inner sep=2pt, font=\footnotesize] (xd) at (0, -0.6) {$x_d$};
    \node[circle, draw=gray!70, fill=gray!10, thick, inner sep=2pt, font=\footnotesize] (xb) at (0, -1.4) {$+1$};

    % Summation Node
    \node[circle, draw=politoNavy, fill=politoNavy!15, very thick, inner sep=5pt, font=\normalsize] (sum) at (2.8, 0.6) {$\sum$};
    \node[font=\tiny\color{politoNavy}, below=2pt of sum] {$z = \mathbf{w}^\top\mathbf{x} + b$};

    % Connections with weights
    \draw[-{Stealth}, thick] (x1) -- node[above, font=\scriptsize] {$w_1$} (sum);
    \draw[-{Stealth}, thick] (x2) -- node[above, font=\scriptsize] {$w_2$} (sum);
    \draw[-{Stealth}, thick] (xd) -- node[below, font=\scriptsize] {$w_d$} (sum);
    \draw[-{Stealth}, thick, dashed] (xb) -- node[below, font=\scriptsize] {$b$} (sum);

    % Activation Function Block
    \node[rectangle, draw=intuitiongreen, fill=intuitiongreen!15, thick, inner sep=5pt, font=\footnotesize, right=0.9cm of sum] (act) {$\operatorname{sgn}(z)$};

    \draw[-{Stealth}, thick] (sum) -- node[above, font=\scriptsize] {$z$} (act);

    % Output
    \node[right=0.9cm of act, font=\footnotesize\bfseries] (out) {$\hat{y} \in \{-1, +1\}$};
    \draw[-{Stealth}, very thick, examred] (act) -- (out);
  \end{scope}
\end{tikzpicture}
\caption{Morphological and mathematical correspondence between biological and artificial neural processing. (a) A biological neuron integrates post-synaptic potentials across dendrites into the metabolic soma, firing an action potential down the myelinated axon toward synaptic terminals. (b) Rosenblatt's artificial perceptron models this mechanism: numerical inputs $x_j$ are modulated by synaptic parameters $w_j$, accumulated by affine summation $z = \mathbf{w}^\top \mathbf{x} + b$, and thresholded by a Heaviside sign function to produce binary classification $\hat{y} \in \{-1, +1\}$.}
\label{fig:biological_vs_artificial_neuron}
\end{figure}

\paragraph{Decision hyperplane geometry.}
Geometrically, the equation $\mathbf{w}^\top \mathbf{x} + b = 0$ defines a $(d-1)$-dimensional \textbf{decision hyperplane} that partitions the $d$-dimensional input space $\mathcal{X} \subseteq \R^d$ into two open half-spaces: a positive half-space where $\mathbf{w}^\top \mathbf{x} + b \ge 0$ (yielding $\hat{y} = +1$) and a negative half-space where $\mathbf{w}^\top \mathbf{x} + b < 0$ (yielding $\hat{y} = -1$). The parameter vector $\mathbf{w} \in \R^d$ is the orthogonal normal vector to the hyperplane pointing directly toward the positive decision half-space. The scalar bias $b \in \R$ offsets the hyperplane from the origin, establishing a signed perpendicular Euclidean distance of $\frac{|b|}{\|\mathbf{w}\|_2}$.

\begin{figure}[H]
\centering
\begin{tikzpicture}[scale=1.15]
  % Axes
  \draw[-{Stealth}, gray!70, thick] (-1.0, 0) -- (5.0, 0) node[right, font=\footnotesize] {$x_1$};
  \draw[-{Stealth}, gray!70, thick] (0, -1.0) -- (0, 4.2) node[above, font=\footnotesize] {$x_2$};

  % Half-space shading
  \fill[deepblue!8, opacity=0.8] (-0.8, 3.8) -- (4.2, 0.46) -- (4.5, 3.8) -- (-0.8, 3.8) -- cycle;
  \fill[examred!8, opacity=0.8] (-0.8, 3.8) -- (4.2, 0.46) -- (4.2, -0.8) -- (-0.8, -0.8) -- cycle;

  % Hyperplane: w1*x1 + w2*x2 + b = 0
  \draw[very thick, politoNavy] (-0.6, 2.5) -- (4.4, -0.8) node[right, font=\footnotesize, align=left] {Decision Hyperplane:\\$\mathbf{w}^\top \mathbf{x} + b = 0$};

  % Normal vector w from a point on the line (1.5, 1.0)
  \coordinate (P) at (1.5, 1.0);
  \draw[-{Stealth[length=6pt]}, ultra thick, intuitiongreen] (P) -- ++(1.0, 1.5) node[above right, font=\footnotesize\bfseries] {$\mathbf{w}$ (Orthogonal Normal Vector)};

  % Perpendicular projection from Origin (0,0) to line
  \coordinate (Proj) at (0.923, 1.385);
  \draw[dashed, thick, gray!80] (0, 0) -- (Proj);
  \draw[thick, gray!80] (0.8, 1.2) -- (0.95, 1.1) -- (1.07, 1.28);
  \node[left, font=\scriptsize\bfseries, text=politoNavy] at (0.45, 0.75) {$d = \frac{|b|}{\|\mathbf{w}\|_2}$};

  % Data points: Positive Class (+1)
  \node[star, star points=5, star point ratio=2.2, draw=deepblue, fill=deepblue!30, thick, inner sep=2.5pt] at (2.2, 2.6) {};
  \node[star, star points=5, star point ratio=2.2, draw=deepblue, fill=deepblue!30, thick, inner sep=2.5pt] at (3.5, 2.1) {};
  \node[star, star points=5, star point ratio=2.2, draw=deepblue, fill=deepblue!30, thick, inner sep=2.5pt] at (1.2, 3.1) {};
  \node[star, star points=5, star point ratio=2.2, draw=deepblue, fill=deepblue!30, thick, inner sep=2.5pt] at (3.8, 3.3) {};
  \node[font=\footnotesize\bfseries, text=deepblue] at (3.2, 3.8) {Positive Half-Space: $\mathbf{w}^\top\mathbf{x} + b > 0$};

  % Data points: Negative Class (-1)
  \node[circle, draw=examred, fill=examred!20, thick, inner sep=2.5pt] at (0.5, 0.5) {};
  \node[circle, draw=examred, fill=examred!20, thick, inner sep=2.5pt] at (1.5, -0.2) {};
  \node[circle, draw=examred, fill=examred!20, thick, inner sep=2.5pt] at (2.8, -0.4) {};
  \node[circle, draw=examred, fill=examred!20, thick, inner sep=2.5pt] at (0.2, 1.2) {};
  \node[font=\footnotesize\bfseries, text=examred] at (1.5, -0.7) {Negative Half-Space: $\mathbf{w}^\top\mathbf{x} + b < 0$};

  % Origin
  \fill (0,0) circle (2pt) node[below left, font=\footnotesize] {$\mathbf{0}$};
\end{tikzpicture}
\caption{The geometry of the linear decision boundary defined by the perceptron. The parameter vector $\mathbf{w} \in \mathbb{R}^d$ is the orthogonal normal vector to the hyperplane pointing directly toward the positive decision half-space. The scalar bias $b \in \mathbb{R}$ offsets the hyperplane from the origin, establishing a signed perpendicular Euclidean distance of $\frac{|b|}{\|\mathbf{w}\|_2}$.}
\label{fig:hyperplane_geometry}
\end{figure}

\paragraph{The perceptron learning algorithm.}
Given a supervised training dataset $\mathcal{D} = \{(\mathbf{x}_i, y_i)\}_{i=1}^n$ with binary labels $y_i \in \{-1, +1\}$, the perceptron learning algorithm operates in an online, mistake-driven fashion. Parameters are initialized arbitrarily (typically $\mathbf{w}^{(0)} = \mathbf{0}, b^{(0)} = 0$).

At step $t$, the algorithm inspects an incoming training instance $(\mathbf{x}_i, y_i)$. If the current hypothesis correctly classifies the instance ($y_i (\mathbf{w}^{(t)\top} \mathbf{x}_i + b^{(t)}) > 0$), no modification is made. If the instance is misclassified ($y_i (\mathbf{w}^{(t)\top} \mathbf{x}_i + b^{(t)}) \le 0$), parameters are updated via the additive correction:
\begin{align}
\mathbf{w}^{(t+1)} &= \mathbf{w}^{(t)} + \eta y_i \mathbf{x}_i, \label{eq:perceptron_update_w} \\
b^{(t+1)} &= b^{(t)} + \eta y_i R^2, \label{eq:perceptron_update_b}
\end{align}
where $\eta > 0$ is the learning rate (often set to $\eta = 1$ without loss of generality, as scaling both sides leaves the sign of predictions unchanged), and $R = \max_i \|\mathbf{x}_i\|_2$ is the radius of the smallest sphere enclosing the data. In parameter space, this update executes vector addition; in input space, it rotates the decision hyperplane directly toward the misclassified sample, pulling it into the correct decision region.

\begin{figure}[H]
\centering
\begin{minipage}{0.48\linewidth}
  \centering
  \begin{tikzpicture}[scale=1.05]
    \node[font=\footnotesize\bfseries\color{politoNavy}] at (1.5, 3.2) {(a) Parameter Space: Vector Correction};
    \draw[-{Stealth}, gray!50] (-0.5, 0) -- (3.2, 0) node[right, font=\tiny] {$w_1$};
    \draw[-{Stealth}, gray!50] (0, -0.5) -- (0, 3.0) node[above, font=\tiny] {$w_2$};
    
    % w(t)
    \coordinate (W0) at (0.8, 1.8);
    \draw[-{Stealth[length=6pt]}, very thick, politoNavy] (0, 0) -- (W0) node[above left, font=\footnotesize] {$\mathbf{w}^{(t)}$};
    
    % Misclassified sample vector y_i * x_i (positive class: +x_i)
    \coordinate (X) at (1.6, -0.4);
    \draw[-{Stealth[length=6pt]}, very thick, examred] (W0) -- ++(X) node[midway, above right, font=\footnotesize] {$+y_i \mathbf{x}_i$};
    
    % w(t+1)
    \coordinate (W1) at ($(W0) + (X)$);
    \draw[-{Stealth[length=6pt]}, very thick, intuitiongreen] (0, 0) -- (W1) node[right, font=\footnotesize\bfseries] {$\mathbf{w}^{(t+1)}$};
    
    \node[draw=politoNavy, fill=politoNavy!5, rounded corners, font=\scriptsize, align=center] at (1.5, -0.8) {$\mathbf{w}^{(t+1)} = \mathbf{w}^{(t)} + \eta y_i \mathbf{x}_i$};
  \end{tikzpicture}
\end{minipage}\hfill
\begin{minipage}{0.48\linewidth}
  \centering
  \begin{tikzpicture}[scale=1.05]
    \node[font=\footnotesize\bfseries\color{politoNavy}] at (1.5, 3.2) {(b) Input Space: Boundary Rotation};
    \draw[-{Stealth}, gray!50] (-1.0, 0) -- (3.2, 0) node[right, font=\tiny] {$x_1$};
    \draw[-{Stealth}, gray!50] (0, -1.0) -- (0, 3.0) node[above, font=\tiny] {$x_2$};

    % Old boundary w(t)^T x = 0
    \draw[thick, politoNavy, dashed] (-0.9, 0.4) -- (2.5, -1.1) node[below right, font=\tiny, align=left] {Old boundary:\\$\mathbf{w}^{(t)\top}\mathbf{x}=0$};
    
    % Misclassified point x_i (positive class y_i = +1)
    \node[star, star points=5, star point ratio=2.2, draw=examred, fill=examred!30, thick, inner sep=2.5pt, label={right:\scriptsize $\mathbf{x}_i \ (y_i = +1)$}] (xi) at (1.2, -0.8) {};

    % New boundary w(t+1)^T x = 0
    \draw[very thick, intuitiongreen] (-0.8, -0.6) -- (2.0, 2.2) node[above right, font=\tiny, align=left] {New boundary:\\$\mathbf{w}^{(t+1)\top}\mathbf{x}=0$};

    % Rotation arrow
    \draw[-{Stealth}, thick, deepdiveTeal] (1.8, -0.7) to[bend right=25] (1.3, 1.4) node[midway, right, font=\tiny] {Boundary rotates toward $\mathbf{x}_i$};
  \end{tikzpicture}
\end{minipage}
\caption{Geometric dynamics of the Perceptron Learning Algorithm update. (a) In parameter space, whenever a training sample $(\mathbf{x}_i, y_i)$ is misclassified, the algorithm performs vector addition $\mathbf{w}^{(t+1)} = \mathbf{w}^{(t)} + y_i \mathbf{x}_i$. (b) In input space, adding $y_i \mathbf{x}_i$ rotates the orthogonal decision hyperplane directly toward the misclassified sample, pulling it into the correct decision region.}
\label{fig:perceptron_learning_dynamics}
\end{figure}

\paragraph{Novikoff's convergence theorem.}
The historical significance of Rosenblatt's perceptron stems from the fact that it was the first machine learning model accompanied by a rigorous mathematical convergence proof, formalized by Albert Novikoff in 1962~\cite{novikoff1962convergence}.

\begin{theorem}[Novikoff's Perceptron Convergence Theorem]
\label{thm:novikoff_convergence}
Let $\mathcal{D} = \{(\mathbf{x}_i, y_i)\}_{i=1}^n$ be a training dataset where $\|\mathbf{x}_i\|_2 \le R$ for all $i$. Assume the dataset is strictly \textbf{linearly separable}: that is, there exists an optimal unit-norm parameter vector $\mathbf{w}^*$ ($\|\mathbf{w}^*\|_2 = 1$) and an optimal margin $\gamma > 0$ such that:
\begin{equation}
y_i (\mathbf{w}^{*\top} \mathbf{x}_i) \ge \gamma, \quad \forall i \in \{1, \dots, n\}.
\label{eq:linear_separability_condition}
\end{equation}
Then, starting from $\mathbf{w}^{(0)} = \mathbf{0}$, the perceptron learning algorithm makes at most $k$ mistake updates before converging to a perfectly separating hyperplane, where:
\begin{equation}
k \le \left( \frac{R}{\gamma} \right)^2.
\label{eq:novikoff_bound}
\end{equation}
\end{theorem}

\begin{proof}
Let $\mathbf{w}^{(k)}$ denote the weight vector after $k$ mistaken updates. Assume the $k$-th update was triggered by misclassified sample $(\mathbf{x}_i, y_i)$, meaning $y_i (\mathbf{w}^{(k-1)\top} \mathbf{x}_i) \le 0$. The update rule is $\mathbf{w}^{(k)} = \mathbf{w}^{(k-1)} + y_i \mathbf{x}_i$.

We bound the inner product $\mathbf{w}^{(k)\top} \mathbf{w}^*$ from below. Expanding recursively:
\begin{equation}
\mathbf{w}^{(k)\top} \mathbf{w}^* = (\mathbf{w}^{(k-1)} + y_i \mathbf{x}_i)^\top \mathbf{w}^* = \mathbf{w}^{(k-1)\top} \mathbf{w}^* + y_i \mathbf{x}_i^\top \mathbf{w}^* \ge \mathbf{w}^{(k-1)\top} \mathbf{w}^* + \gamma.
\end{equation}
Inductively from $\mathbf{w}^{(0)} = \mathbf{0}$, after $k$ mistakes:
\begin{equation}
\mathbf{w}^{(k)\top} \mathbf{w}^* \ge k \gamma.
\label{eq:novikoff_inner_bound}
\end{equation}

Next, we bound the squared Euclidean norm $\|\mathbf{w}^{(k)}\|_2^2$ from above:
\begin{align}
\|\mathbf{w}^{(k)}\|_2^2 &= \|\mathbf{w}^{(k-1)} + y_i \mathbf{x}_i\|_2^2 = \|\mathbf{w}^{(k-1)}\|_2^2 + 2 y_i \mathbf{w}^{(k-1)\top} \mathbf{x}_i + \|\mathbf{x}_i\|_2^2 \nonumber \\
&\le \|\mathbf{w}^{(k-1)}\|_2^2 + 0 + R^2 = \|\mathbf{w}^{(k-1)}\|_2^2 + R^2,
\end{align}
where we used the mistake condition $y_i \mathbf{w}^{(k-1)\top} \mathbf{x}_i \le 0$ and the radius bound $\|\mathbf{x}_i\|_2 \le R$. Inductively from $\mathbf{w}^{(0)} = \mathbf{0}$:
\begin{equation}
\|\mathbf{w}^{(k)}\|_2^2 \le k R^2 \implies \|\mathbf{w}^{(k)}\|_2 \le \sqrt{k} R.
\label{eq:novikoff_norm_bound}
\end{equation}

By the Cauchy-Schwarz inequality, for any unit vector $\mathbf{w}^*$:
\begin{equation}
\mathbf{w}^{(k)\top} \mathbf{w}^* \le \|\mathbf{w}^{(k)}\|_2 \|\mathbf{w}^*\|_2 = \|\mathbf{w}^{(k)}\|_2.
\end{equation}
Combining Equation~\eqref{eq:novikoff_inner_bound} and Equation~\eqref{eq:novikoff_norm_bound}:
\begin{equation}
k \gamma \le \mathbf{w}^{(k)\top} \mathbf{w}^* \le \|\mathbf{w}^{(k)}\|_2 \le \sqrt{k} R \implies k \gamma \le \sqrt{k} R \implies \sqrt{k} \le \frac{R}{\gamma} \implies k \le \left( \frac{R}{\gamma} \right)^2.
\end{equation}
This completes the proof.
\end{proof}

\subsection{The fatal flaws}
\label{sec:xor_and_minsky_critique}

Despite the mathematical elegance of Novikoff's theorem, the single-layer perceptron suffers from fundamental theoretical limitations that ultimately triggered the historical First AI Winter.

\paragraph{The non-separability catastrophe: Boundary thrashing and infinite cycling.}
Novikoff's Perceptron Convergence Theorem (Theorem~\ref{thm:novikoff_convergence}) holds strictly under the premise that the training dataset $\mathcal{D}$ is linearly separable, meaning that a strictly positive margin exists: $\gamma > 0$.

However, in realistic real-world datasets, sensory noise, label ambiguity, and overlapping class probability distributions inevitably render data \textbf{linearly non-separable} ($\gamma \le 0$). Under non-separability, the convergence upper bound $k \le (R/\gamma)^2$ becomes undefined ($\gamma \to 0$ implies $k \to \infty$).

Crucially, the presence of an inseparable outlier or overlapping class instance (such as a negative point located within a cluster of positive samples) triggers a catastrophic failure mode:
\begin{itemize}[noitemsep,topsep=2pt,parsep=0pt,partopsep=0pt]
  \item \textbf{Boundary thrashing and non-terminating ping-pong oscillation}: Because the outlier violates the margin condition, the algorithm executes the additive update $\mathbf{w}^{(t+1)} = \mathbf{w}^{(t)} + y_i \mathbf{x}_i$, forcibly rotating the decision boundary from $H^{(t)}$ to $H^{(t+1)}$ to satisfy that outlier. However, this corrective rotation inevitably sweeps across previously well-classified instances, instantly converting them into misclassified errors. In subsequent iterations, these newly displaced points trigger counter-updates that swing the boundary back toward $H^{(t)}$. The algorithm thus becomes trapped in an infinite ping-pong loop, thrashing perpetually between conflicting hyperplanes without ever terminating ($k \to \infty$).
\end{itemize}

To mitigate this catastrophic behavior in practical applications, Stephen Gallant introduced the heuristic \textbf{Pocket Algorithm}~\cite{gallant1990perceptron}, which runs the perceptron algorithm for a predetermined number of iterations while storing (``in its pocket'') the parameter vector that achieved the lowest classification error seen so far. Nevertheless, the fundamental theoretical guarantee of exact convergence is irrevocably lost on non-separable data.

\begin{figure}[H]
\centering
\begin{tikzpicture}[scale=1.15]
  \draw[-{Stealth}, gray!50] (-0.5, 0) -- (4.4, 0) node[right, font=\footnotesize] {$x_1$};
  \draw[-{Stealth}, gray!50] (0, -0.5) -- (0, 3.6) node[above, font=\footnotesize] {$x_2$};

  % Positive cluster (+1 stars)
  \node[star, star points=5, star point ratio=2.2, draw=politoNavy, fill=politoNavy!30, thick, inner sep=2pt] at (0.8, 2.5) {};
  \node[star, star points=5, star point ratio=2.2, draw=politoNavy, fill=politoNavy!30, thick, inner sep=2pt] at (1.3, 3.1) {};
  \node[star, star points=5, star point ratio=2.2, draw=politoNavy, fill=politoNavy!30, thick, inner sep=2pt] at (2.2, 3.0) {};
  \node[star, star points=5, star point ratio=2.2, draw=politoNavy, fill=politoNavy!30, thick, inner sep=2pt] at (2.5, 2.4) {};

  % Negative cluster (-1 circles)
  \node[circle, draw=examred, fill=examred!20, thick, inner sep=2pt] at (2.7, 0.7) {};
  \node[circle, draw=examred, fill=examred!20, thick, inner sep=2pt] at (3.4, 1.1) {};
  \node[circle, draw=examred, fill=examred!20, thick, inner sep=2pt] at (3.2, 0.4) {};

  % The Inseparable Outlier Negative point inside Positive cluster
  \node[circle, draw=examred, fill=examred!70, very thick, inner sep=2.5pt] (outlier) at (1.5, 2.0) {};
  \node[font=\scriptsize\bfseries, text=examred, above left] at (1.5, 2.0) {Inseparable Outlier ($-1$)};

  % Two alternating oscillating hyperplanes
  \draw[thick, deepblue, dashed] (0.2, 1.4) -- (3.8, 2.6) node[right, font=\scriptsize, text=deepblue] {Hyperplane $H^{(t)}$};
  \draw[thick, examred, dashdotted] (0.8, 3.4) -- (2.6, 0.3) node[below right, font=\scriptsize, text=examred] {Hyperplane $H^{(t+1)}$};

  % Oscillation cycle arrows in clear zone
  \draw[-{Stealth}, thick, deepdiveTeal] (3.3, 2.2) to[bend left=35] (2.9, 1.2);
  \draw[-{Stealth}, thick, deepdiveTeal] (2.7, 1.2) to[bend left=35] (3.1, 2.2);
  \node[font=\scriptsize\bfseries, text=deepdiveTeal, right, align=left] at (3.4, 1.7) {Perpetual\\Oscillation};

  \node[draw=examred, fill=examback, rounded corners, font=\scriptsize, align=center] at (2.0, -0.6) {Non-separable data: Algorithm oscillates indefinitely ($k \to \infty$).};
\end{tikzpicture}
\caption{Flaw 1 of the Perceptron: Non-separability and perpetual boundary thrashing. When classes overlap or contain an inseparable outlier, no separating hyperplane exists ($\gamma \le 0$). The perceptron learning algorithm continually rotates to satisfy the misclassified outlier, which promptly causes previously classified points to fail, trapping the algorithm in an infinite ping-pong oscillation between conflicting hyperplanes $H^{(t)}$ and $H^{(t+1)}$.}
\label{fig:perceptron_flaw_nonseparable}
\end{figure}

\paragraph{Inability to compute Exclusive-OR (XOR).}
The non-separability catastrophe demonstrates that the perceptron fails in the presence of noise and outliers. Yet a far deeper question remains: what if the data contains zero noise, zero outliers, and is completely deterministic? Can a single-layer perceptron learn even the simplest logical functions?

In their landmark 1969 monograph \textit{Perceptrons}, Marvin Minsky and Seymour Papert~\cite{minsky1969perceptrons} established that a single-layer perceptron is geometrically incapable of computing basic non-linear Boolean predicates, most famously the \textbf{Exclusive-OR (XOR)} function.

Consider the four canonical vertices of the unit square $[0, 1]^2$ defining the XOR logic gate:
\begin{equation}
\begin{aligned}
\mathbf{x}_1 = [0, 0]^\top &\implies y = -1, \\
\mathbf{x}_2 = [0, 1]^\top &\implies y = +1,
\end{aligned}
\qquad
\begin{aligned}
\mathbf{x}_3 = [1, 0]^\top &\implies y = +1, \\
\mathbf{x}_4 = [1, 1]^\top &\implies y = -1.
\end{aligned}
\label{eq:xor_dataset}
\end{equation}

To classify these four points correctly with zero training error, a perceptron parameterized by weights $(w_1, w_2)$ and bias $b$ must satisfy the four simultaneous strict linear inequalities:
\begin{align}
\text{For } \mathbf{x}_1 = [0, 0]^\top: \quad & w_1 \cdot 0 + w_2 \cdot 0 + b < 0 \implies b < 0, \label{eq:xor_ineq1} \\
\text{For } \mathbf{x}_2 = [0, 1]^\top: \quad & w_1 \cdot 0 + w_2 \cdot 1 + b \ge 0 \implies w_2 + b \ge 0, \label{eq:xor_ineq2} \\
\text{For } \mathbf{x}_3 = [1, 0]^\top: \quad & w_1 \cdot 1 + w_2 \cdot 0 + b \ge 0 \implies w_1 + b \ge 0, \label{eq:xor_ineq3} \\
\text{For } \mathbf{x}_4 = [1, 1]^\top: \quad & w_1 \cdot 1 + w_2 \cdot 1 + b < 0 \implies w_1 + w_2 + b < 0. \label{eq:xor_ineq4}
\end{align}

We now derive the formal algebraic impossibility proof:
\begin{proof}
Summing the two positive-class inequalities~\eqref{eq:xor_ineq2} and~\eqref{eq:xor_ineq3}:
\begin{equation}
(w_2 + b) + (w_1 + b) \ge 0 \implies w_1 + w_2 + 2b \ge 0 \implies w_1 + w_2 + b \ge -b.
\end{equation}
Because Equation~\eqref{eq:xor_ineq1} requires $b < 0$, multiplying by $-1$ yields $-b > 0$. Therefore:
\begin{equation}
w_1 + w_2 + b \ge -b > 0 \implies w_1 + w_2 + b > 0.
\label{eq:xor_contradiction_positive}
\end{equation}
However, Equation~\eqref{eq:xor_ineq4} strictly requires:
\begin{equation}
w_1 + w_2 + b < 0.
\label{eq:xor_contradiction_negative}
\end{equation}
Equations~\eqref{eq:xor_contradiction_positive} and~\eqref{eq:xor_contradiction_negative} are directly contradictory: a real quantity cannot be strictly positive and strictly negative simultaneously ($0 < w_1 + w_2 + b < 0$, a manifest contradiction in $\mathbb{R}$).
\end{proof}

Hence, \textbf{no parameter vector $\mathbf{w} \in \mathbb{R}^2$ and bias $b \in \mathbb{R}$ can ever classify the XOR problem}. Geometrically, as illustrated in Figure~\ref{fig:perceptron_flaw_xor}, the positive points $(0, 1)$ and $(1, 0)$ and negative points $(0, 0)$ and $(1, 1)$ occupy opposite diagonals of the unit square. Any straight line drawn through the plane can at best isolate three of the four points, always misclassifying at least one vertex.

\begin{figure}[H]
\centering
\begin{tikzpicture}[scale=1.3]
  % Axes
  \draw[-{Stealth}, gray!70, thick] (-0.4, 0) -- (2.2, 0) node[right, font=\footnotesize] {$x_1$};
  \draw[-{Stealth}, gray!70, thick] (0, -0.4) -- (0, 2.2) node[above, font=\footnotesize] {$x_2$};
  \draw[help lines, color=gray!20, step=1] (-0.2,-0.2) grid (1.5,1.5);

  % XOR points with non-conflicting placement
  \node[circle, draw=examred, fill=examred!30, very thick, inner sep=3pt, label={below left:\footnotesize $(0,0): y=-1$}] at (0, 0) {};
  \node[circle, draw=examred, fill=examred!30, very thick, inner sep=3pt, label={above right:\footnotesize $(1,1): y=-1$}] at (1, 1) {};
  \node[star, star points=5, star point ratio=2.2, draw=politoNavy, fill=politoNavy!40, very thick, inner sep=3pt, label={above left:\footnotesize $(0,1): y=+1$}] at (0, 1) {};
  \node[star, star points=5, star point ratio=2.2, draw=politoNavy, fill=politoNavy!40, very thick, inner sep=3pt, label={below right:\footnotesize $(1,0): y=+1$}] at (1, 0) {};

  % Candidate hyperplanes failing (placed in open spaces)
  \draw[thick, gray!60, dashed] (-0.3, 0.5) -- (1.7, 0.5) node[right, font=\tiny, text=gray!80] {Candidate 1 ($x_2 = 0.5$)};
  \draw[thick, gray!60, dashed] (0.5, -0.3) -- (0.5, 1.7) node[above, font=\tiny, text=gray!80] {Candidate 2 ($x_1 = 0.5$)};
  \draw[thick, examred, dashdotted] (-0.3, 0.8) node[left, font=\tiny, text=examred] {Diagonal partition} -- (1.1, -0.2);

  \node[draw=examred, fill=examback, rounded corners, font=\scriptsize, align=center] at (0.9, -0.9) {Minsky \& Papert (1969): No single linear boundary can isolate opposite diagonal pairs.};
\end{tikzpicture}
\caption{Flaw 2 of the Perceptron: The XOR non-linear separability impasse. The Exclusive-OR (XOR) function assigns positive labels ($y = +1$) to coordinates $(0, 1)$ and $(1, 0)$, and negative labels ($y = -1$) to $(0, 0)$ and $(1, 1)$. Because positive and negative pairs occupy alternate diagonals of the unit square $[0, 1]^2$, no single linear decision hyperplane can separate the two classes.}
\label{fig:perceptron_flaw_xor}
\end{figure}

The publication of Minsky and Papert's critique had devastating historical ramifications. Because early AI researchers lacked an algorithm to train multi-layer neural architectures with hidden units, the field broadly concluded that artificial neural networks were fundamentally flawed and unscalable, causing government and academic funding to evaporate and initiating the decade-long First AI Winter.

\paragraph{Zero derivative and optimization incompatibility.}
The third fatal flaw of Rosenblatt's perceptron lies in its activation function. The discontinuous step activation:
\begin{equation}
f(z) = \operatorname{sgn}(z) = \begin{cases} +1 & \text{if } z \ge 0, \\ -1 & \text{if } z < 0, \end{cases}
\end{equation}
is piecewise constant. Its derivative is zero everywhere that it is defined:
\begin{equation}
f'(z) = \frac{d}{dz}\operatorname{sgn}(z) = 0, \quad \forall z \ne 0,
\end{equation}
and is completely undefined (or infinite) at the discontinuity $z = 0$. 

Consequently, the parameter gradient of the pre-activation with respect to weights:
\begin{equation}
\nabla_\mathbf{w} f(\mathbf{w}^\top \mathbf{x} + b) = f'(z) \cdot \mathbf{x} = \mathbf{0}, \quad \forall z \ne 0.
\end{equation}
Because the gradient is identically zero almost everywhere, \textbf{gradient-based optimization algorithms cannot receive any directional learning signal}. Perturbing the weights by an infinitesimal step $\Delta \mathbf{w}$ either produces zero change in output or an abrupt, unpredictable jump across the step discontinuity. 

The perceptron's inability to provide smooth gradient signals, coupled with its linear separability barrier on problems like XOR, established the definitive conclusion of the classical era: single-layer linear architectures cannot serve as universal intelligent learners. To break through this barrier, machine learning had to transition to layered compositions of continuous, differentiable units---leading directly to the Multi-Layer Perceptron in Section~\ref{sec:mlp_and_representation_learning}.

\section{Multilayer perceptrons and approximation}
\label{sec:mlp_and_representation_learning}

To liberate neural computing from the linear separability impasse, the architecture of the artificial neuron must be fundamentally reconceived. The failure on XOR and the non-differentiability of step functions revealed that a single hyperplane cutting through raw input space cannot capture complex decision geometries. A learner must be capable of transforming the geometry of the input space itself.

However, resolving this challenge is not as straightforward as simply stacking multiple linear perceptrons into a cascading sequence. As linear algebra dictates, any sequence of cascaded affine transformations collapses into a single affine transformation, offering zero representational advantage over a single-layer model. To unlock true hierarchical expressiveness, two simultaneous transformations are mathematically mandatory: (1) interleaving linear matrix projections with continuous, non-linear activation functions that break affine subspace closure, and (2) organizing neurons into intermediate \textbf{hidden layers} that execute automated \textbf{representation learning}.

\subsection{Beyond linearity and representation}
\label{sec:continuous_probabilities_softmax}
\label{sec:resolving_xor_hidden_representations}

\paragraph{Why we must go beyond linearity.}
In classical linear models, decision boundaries are strictly constrained to $(D-1)$-dimensional affine hyperplanes of the form $\mathbf{w}^\top \mathbf{x} + b = 0$. While computationally tractable and admitting closed-form solutions in regression tasks, linear decision boundaries are fundamentally incapable of modeling complex real-world data distributions. Perceptual data---such as natural images, acoustic speech signals, and semantic text---invariably lie on highly curved, non-linear sub-manifolds embedded within high-dimensional ambient space. 

When class distributions exhibit topological complexity, non-convex geometry, or inter-class interleaving---as exemplified by the classic Exclusive-OR (XOR) logic function---a linear boundary cannot isolate the classes without incurring catastrophic classification error. To build expressive learning machines, we must endow models with the mathematical capacity to construct non-linear decision boundaries.

\paragraph{The compositional collapse of purely linear networks.}
A naive intuition might suggest that multi-layer expressiveness could be attained simply by cascading multiple linear layers, feeding the output of one linear perceptron as input to the next. Consider an $L$-layer network where each intermediate layer computes an affine transformation without activation functions:
\begin{align}
\mathbf{h}^{(1)} &= \mathbf{W}_1 \mathbf{x} + \mathbf{b}_1, \nonumber \\
\mathbf{h}^{(2)} &= \mathbf{W}_2 \mathbf{h}^{(1)} + \mathbf{b}_2 = \mathbf{W}_2 (\mathbf{W}_1 \mathbf{x} + \mathbf{b}_1) + \mathbf{b}_2 = (\mathbf{W}_2 \mathbf{W}_1)\mathbf{x} + (\mathbf{W}_2 \mathbf{b}_1 + \mathbf{b}_2), \nonumber \\
&\;\;\vdots \nonumber \\
\hat{\mathbf{y}} &= \mathbf{W}_L \mathbf{h}^{(L-1)} + \mathbf{b}_L.
\label{eq:linear_network_cascade}
\end{align}

\begin{proof}
We prove by mathematical induction on the layer depth $L$ that any cascading sequence of affine maps collapses to a single affine transformation.

\textit{Base Case ($L=1$):} For a single layer, $\hat{\mathbf{y}} = \mathbf{W}_1 \mathbf{x} + \mathbf{b}_1$, which is trivially affine with effective weight matrix $\mathbf{W}_{\text{eff}}^{(1)} = \mathbf{W}_1$ and bias $\mathbf{b}_{\text{eff}}^{(1)} = \mathbf{b}_1$.

\textit{Inductive Hypothesis:} Assume that for any layer $k \ge 1$, the representation $\mathbf{h}^{(k)}$ is an affine function of the input $\mathbf{x}$:
\begin{equation}
\mathbf{h}^{(k)} = \mathbf{W}_{\text{eff}}^{(k)} \mathbf{x} + \mathbf{b}_{\text{eff}}^{(k)}, \quad \text{with } \mathbf{W}_{\text{eff}}^{(k)} \in \R^{d_k \times d_0}, \; \mathbf{b}_{\text{eff}}^{(k)} \in \R^{d_k}.
\end{equation}

\textit{Inductive Step:} Applying the affine transformation of layer $k+1$ yields:
\begin{align}
\mathbf{h}^{(k+1)} &= \mathbf{W}_{k+1} \mathbf{h}^{(k)} + \mathbf{b}_{k+1} \nonumber \\
&= \mathbf{W}_{k+1} \left( \mathbf{W}_{\text{eff}}^{(k)} \mathbf{x} + \mathbf{b}_{\text{eff}}^{(k)} \right) + \mathbf{b}_{k+1} \nonumber \\
&= \left( \mathbf{W}_{k+1} \mathbf{W}_{\text{eff}}^{(k)} \right) \mathbf{x} + \left( \mathbf{W}_{k+1} \mathbf{b}_{\text{eff}}^{(k)} + \mathbf{b}_{k+1} \right).
\end{align}
Defining the composite parameter matrices:
\begin{equation}
\mathbf{W}_{\text{eff}}^{(k+1)} \triangleq \mathbf{W}_{k+1} \mathbf{W}_{\text{eff}}^{(k)} \in \R^{d_{k+1} \times d_0}, \quad \mathbf{b}_{\text{eff}}^{(k+1)} \triangleq \mathbf{W}_{k+1} \mathbf{b}_{\text{eff}}^{(k)} + \mathbf{b}_{k+1} \in \R^{d_{k+1}},
\end{equation}
the representation $\mathbf{h}^{(k+1)}$ remains strictly an affine function of $\mathbf{x}$. By induction, an $L$-layer purely linear network computes:
\begin{equation}
\hat{\mathbf{y}} = \mathbf{W}_{\text{eff}} \mathbf{x} + \mathbf{b}_{\text{eff}}, \quad \text{with } \mathbf{W}_{\text{eff}} = \prod_{l=L}^1 \mathbf{W}_l = \mathbf{W}_L \mathbf{W}_{L-1} \cdots \mathbf{W}_1,
\label{eq:compositional_collapse_result}
\end{equation}
and $\mathbf{b}_{\text{eff}} = \mathbf{b}_L + \sum_{l=1}^{L-1} \left( \prod_{j=L}^{l+1} \mathbf{W}_j \right) \mathbf{b}_l$.
\end{proof}

Because matrix multiplication is associative and affine transformations form a mathematically closed set under functional composition, \textbf{stacking arbitrary depths of purely linear layers yields zero additional representational power over a single-layer linear model}. No matter how many millions of linear units are assembled, the composition remains incapable of bending a hyperplane.

\paragraph{The imperative for non-linear activation functions.}
To escape compositional collapse, we must break affine subspace closure. The mathematical remedy is to interleave each linear matrix projection with an element-wise, non-linear activation function $\sigma: \R \to \R$:
\begin{equation}
\mathbf{h}^{(l)} = \sigma\left( \mathbf{W}_l \mathbf{h}^{(l-1)} + \mathbf{b}_l \right), \quad \forall l \in \{1, \dots, L\}.
\label{eq:interleaved_nonlinear_layer}
\end{equation}
Because $\sigma(\cdot)$ is non-linear, the operator $\sigma$ cannot commute with matrix multiplication: $\mathbf{W}_2 \sigma(\mathbf{W}_1 \mathbf{x} + \mathbf{b}_1) \ne \tilde{\mathbf{W}} \mathbf{x} + \tilde{\mathbf{b}}$. The non-linearity prevents the successive weight matrices from collapsing into a single matrix product. Geometrically, non-linear activations fold, stretch, and twist the coordinate axes, allowing the network to carve curved, multi-region decision boundaries that solve topologically complex tasks.

\paragraph{Resolving the XOR dilemma via representation learning.}
To see concretely how non-linear hidden representations shatter the linear separability barrier, consider Minsky and Papert's fatal XOR counterexample. The XOR function assigns positive labels ($y=1$) to $[1, 0]^\top$ and $[0, 1]^\top$, and negative labels ($y=0$) to $[0, 0]^\top$ and $[1, 1]^\top$.

We construct a minimal two-layer feedforward network with two input units, two hidden units equipped with the Rectified Linear Unit ($\operatorname{ReLU}(z) \triangleq \max(0, z)$), and one linear output unit:
\begin{align}
\mathbf{h} &= \operatorname{ReLU}(\mathbf{W}_1 \mathbf{x} + \mathbf{b}_1), \quad \text{with } \mathbf{W}_1 = \begin{bmatrix} 1 & 1 \\ 1 & 1 \end{bmatrix}, \, \mathbf{b}_1 = \begin{bmatrix} 0 \\ -1 \end{bmatrix}, \label{eq:xor_hidden_forward} \\
\hat{y} &= \mathbf{w}_2^\top \mathbf{h} + b_2, \quad \text{with } \mathbf{w}_2 = \begin{bmatrix} 1 \\ -2 \end{bmatrix}, \, b_2 = 0. \label{eq:xor_output_forward}
\end{align}

\begin{figure}[H]
\centering
\begin{tikzpicture}[scale=1.0, every node/.style={transform shape}]
  % Layer 0: Input Layer
  \node[circle, draw=deepblue, fill=deepblue!10, thick, inner sep=3pt] (x1) at (0, 1.2) {$x_1$};
  \node[circle, draw=deepblue, fill=deepblue!10, thick, inner sep=3pt] (x2) at (0, -1.2) {$x_2$};
  \node[font=\footnotesize\bfseries\color{politoNavy}] at (0, 2.2) {Input $\mathbf{x} \in \R^2$};

  % Layer 1: Hidden Layer
  \node[circle, draw=intuitiongreen, fill=intuitiongreen!15, very thick, inner sep=3pt] (h1) at (3.2, 1.2) {$h_1$};
  \node[circle, draw=intuitiongreen, fill=intuitiongreen!15, very thick, inner sep=3pt] (h2) at (3.2, -1.2) {$h_2$};
  \node[font=\footnotesize\bfseries\color{intuitiongreen}] at (3.2, 2.2) {Hidden $\mathbf{h} = \operatorname{ReLU}(\mathbf{W}_1\mathbf{x}+\mathbf{b}_1)$};

  % Layer 2: Output Layer
  \node[circle, draw=examred, fill=examred!15, very thick, inner sep=3pt] (y) at (6.4, 0.0) {$\hat{y}$};
  \node[font=\footnotesize\bfseries\color{examred}] at (6.4, 2.2) {Output $\hat{y} = \mathbf{w}_2^\top\mathbf{h}+b_2$};

  % Connections: Input to Hidden (W1)
  \draw[-{Stealth}, thick] (x1) -- node[above, font=\scriptsize] {$W_{11}=1$} (h1);
  \draw[-{Stealth}, thick] (x1) -- node[near start, above right, font=\scriptsize] {$W_{21}=1$} (h2);
  \draw[-{Stealth}, thick] (x2) -- node[near start, below right, font=\scriptsize] {$W_{12}=1$} (h1);
  \draw[-{Stealth}, thick] (x2) -- node[below, font=\scriptsize] {$W_{22}=1$} (h2);

  % Biases for Hidden
  \node[gray!70, font=\scriptsize] (b1_1) at (3.2, 0.4) {$b_1=0$};
  \node[gray!70, font=\scriptsize] (b1_2) at (3.2, -2.0) {$b_2=-1$};

  % Connections: Hidden to Output (w2)
  \draw[-{Stealth}, thick] (h1) -- node[above right, font=\scriptsize] {$w_{2,1}=1$} (y);
  \draw[-{Stealth}, thick] (h2) -- node[below right, font=\scriptsize] {$w_{2,2}=-2$} (y);

  % Output Arrow
  \draw[-{Stealth}, very thick, examred] (y) -- ++(1.2, 0) node[right, font=\footnotesize\bfseries] {$\hat{y} \in \{0, 1\}$};
\end{tikzpicture}
\caption{The Multi-Layer Perceptron (MLP) architecture resolving XOR. Input coordinates $(x_1, x_2)$ are projected into a 2D hidden layer $\mathbf{h} = \operatorname{ReLU}(\mathbf{W}_1 \mathbf{x} + \mathbf{b}_1)$ that extracts non-linear representations, followed by a linear output unit that classifies the transformed points cleanly.}
\label{fig:mlp_architecture}
\end{figure}

Evaluating the forward pass step-by-step across all four boolean input combinations:
\begin{enumerate}[noitemsep,topsep=2pt,parsep=0pt,partopsep=0pt]
  \item \textbf{Input $\mathbf{x} = [0, 0]^\top$ (Class 0)}:
    \begin{equation}
    \mathbf{W}_1 \mathbf{x} + \mathbf{b}_1 = \begin{bmatrix} 0 \\ -1 \end{bmatrix} \implies \mathbf{h} = \operatorname{ReLU}\left(\begin{bmatrix} 0 \\ -1 \end{bmatrix}\right) = \begin{bmatrix} 0 \\ 0 \end{bmatrix} \implies \hat{y} = [1, -2] \begin{bmatrix} 0 \\ 0 \end{bmatrix} = 0.
    \end{equation}
  \item \textbf{Input $\mathbf{x} = [1, 0]^\top$ (Class 1)}:
    \begin{equation}
    \mathbf{W}_1 \mathbf{x} + \mathbf{b}_1 = \begin{bmatrix} 1 \\ 0 \end{bmatrix} \implies \mathbf{h} = \operatorname{ReLU}\left(\begin{bmatrix} 1 \\ 0 \end{bmatrix}\right) = \begin{bmatrix} 1 \\ 0 \end{bmatrix} \implies \hat{y} = [1, -2] \begin{bmatrix} 1 \\ 0 \end{bmatrix} = 1.
    \end{equation}
  \item \textbf{Input $\mathbf{x} = [0, 1]^\top$ (Class 1)}:
    \begin{equation}
    \mathbf{W}_1 \mathbf{x} + \mathbf{b}_1 = \begin{bmatrix} 1 \\ 0 \end{bmatrix} \implies \mathbf{h} = \operatorname{ReLU}\left(\begin{bmatrix} 1 \\ 0 \end{bmatrix}\right) = \begin{bmatrix} 1 \\ 0 \end{bmatrix} \implies \hat{y} = [1, -2] \begin{bmatrix} 1 \\ 0 \end{bmatrix} = 1.
    \end{equation}
  \item \textbf{Input $\mathbf{x} = [1, 1]^\top$ (Class 0)}:
    \begin{equation}
    \mathbf{W}_1 \mathbf{x} + \mathbf{b}_1 = \begin{bmatrix} 2 \\ 1 \end{bmatrix} \implies \mathbf{h} = \operatorname{ReLU}\left(\begin{bmatrix} 2 \\ 1 \end{bmatrix}\right) = \begin{bmatrix} 2 \\ 1 \end{bmatrix} \implies \hat{y} = [1, -2] \begin{bmatrix} 2 \\ 1 \end{bmatrix} = 2 - 2 = 0.
    \end{equation}
\end{enumerate}
Every input point is classified with zero error ($\hat{y} = y$).

The geometric explanation behind this triumph is illustrated in Figure~\ref{fig:xor_representation_warping}. In the original input space $\mathcal{X}$, the positive instances $[1, 0]^\top$ and $[0, 1]^\top$ are positioned along the off-diagonal, diagonally interleaved between the negative instances $[0, 0]^\top$ and $[1, 1]^\top$; no straight line can separate them. The non-linear hidden transformation $\mathbf{h} = \operatorname{ReLU}(\mathbf{W}_1 \mathbf{x} + \mathbf{b}_1)$ acts as a coordinate folding operation: it projects the two separate positive points $[1, 0]^\top$ and $[0, 1]^\top$ onto the exact same coordinate $[1, 0]^\top$ in the latent hidden space $\mathcal{H}$, while mapping $[0, 0]^\top \to [0, 0]^\top$ and $[1, 1]^\top \to [2, 1]^\top$. In this newly formed hidden representation space $\mathcal{H}$, the class distributions are strictly linearly separable by the hyperplane $h_1 - 2h_2 = 0.5$.

\begin{figure}[H]
\centering
\begin{minipage}{0.48\linewidth}
  \centering
  \begin{tikzpicture}[scale=1.1]
    \draw[-{Stealth}, gray!70, thick] (-0.3, 0) -- (1.6, 0) node[right, font=\footnotesize] {$x_1$};
    \draw[-{Stealth}, gray!70, thick] (0, -0.3) -- (0, 1.6) node[above, font=\footnotesize] {$x_2$};
    \draw[help lines, color=gray!20, step=1] (-0.2,-0.2) grid (1.4,1.4);
    
    % Points
    \node[circle, draw=examred, fill=examred!20, thick, inner sep=3pt, label={below left:\scriptsize $(0,0): y=0$}] at (0, 0) {};
    \node[circle, draw=examred, fill=examred!20, thick, inner sep=3pt, label={above right:\scriptsize $(1,1): y=0$}] at (1, 1) {};
    \node[star, star points=5, star point ratio=2.2, draw=politoNavy, fill=politoNavy!30, thick, inner sep=2.5pt, label={above left:\scriptsize $(0,1): y=1$}] at (0, 1) {};
    \node[star, star points=5, star point ratio=2.2, draw=politoNavy, fill=politoNavy!30, thick, inner sep=2.5pt, label={below right:\scriptsize $(1,0): y=1$}] at (1, 0) {};

    \node[font=\footnotesize\bfseries, align=center, text=examred] at (0.5, -0.7) {Input Space $\mathcal{X}$:\\Non-linearly Separable};
  \end{tikzpicture}
\end{minipage}\hfill
\begin{minipage}{0.48\linewidth}
  \centering
  \begin{tikzpicture}[scale=1.1]
    \draw[-{Stealth}, gray!70, thick] (-0.3, 0) -- (2.4, 0) node[right, font=\footnotesize] {$h_1$};
    \draw[-{Stealth}, gray!70, thick] (0, -0.3) -- (0, 1.6) node[above, font=\footnotesize] {$h_2$};
    \draw[help lines, color=gray!20, step=1] (-0.2,-0.2) grid (2.2,1.4);

    % Transformed Points
    \node[circle, draw=examred, fill=examred!20, thick, inner sep=3pt, label={below left:\scriptsize $(0,0)$}] at (0, 0) {};
    \node[circle, draw=examred, fill=examred!20, thick, inner sep=3pt, label={above right:\scriptsize $(2,1)$}] at (2, 1) {};
    \node[star, star points=5, star point ratio=2.2, draw=politoNavy, fill=politoNavy!30, thick, inner sep=2.5pt, label={above:\scriptsize Coincident: $(1,0)$}] at (1, 0) {};

    % Separating hyperplane
    \draw[very thick, intuitiongreen] (0.5, -0.2) -- (1.5, 1.4) node[right, font=\scriptsize, align=left] {Hyperplane:\\$h_1 - 2h_2 = 0.5$};

    \node[font=\footnotesize\bfseries, align=center, text=intuitiongreen] at (1.0, -0.7) {Hidden Space $\mathcal{H}$:\\Linearly Separable!};
  \end{tikzpicture}
\end{minipage}
\caption{The geometric mechanism of representation learning resolving the XOR problem. Left: In the raw input space $\mathcal{X}$, no single line can separate the positive stars from the negative circles. Right: The non-linear hidden layer transformation $\mathbf{h} = \operatorname{ReLU}(\mathbf{W}_1 \mathbf{x} + \mathbf{b}_1)$ collapses both positive points $[0, 1]^\top$ and $[1, 0]^\top$ onto the single coordinate $[1, 0]^\top$, while mapping $[1, 1]^\top$ to $[2, 1]^\top$. In this hidden space $\mathcal{H}$, the classes are cleanly separated by a single linear hyperplane $h_1 - 2h_2 = 0.5$.}
\label{fig:xor_representation_warping}
\end{figure}

This reveals the foundational paradigm of modern deep learning: \textbf{representation learning}. Rather than requiring human engineers to manually design domain-specific features to make inputs linearly separable, the hidden layer of a neural network autonomously transforms raw coordinate space into a latent geometry where linear classification succeeds.

\subsection{Universal approximation and depth}
\label{sec:universal_approximation_theorem}
\label{sec:perceptron_to_fcn}
\label{sec:activation_functions_foundations}

\paragraph{The Universal Approximation Theorem.}
How expressive is a feedforward neural network equipped with non-linear activation functions? The foundational theoretical answer is established by the celebrated \textbf{Universal Approximation Theorem}, proven by George Cybenko (1989)~\cite{cybenko1989approximation} for sigmoidal activations and generalized by Kurt Hornik (1991)~\cite{hornik1991approximation} to arbitrary continuous, non-constant activations (and subsequently extended to non-polynomial continuous activations including ReLU).

\begin{theorem}[Universal Approximation Theorem~\cite{cybenko1989approximation,hornik1991approximation}]
\label{thm:uat}
Let $\sigma: \R \to \R$ be any continuous, non-constant, bounded (or continuous non-polynomial) activation function. Let $K \subset \R^D$ be a compact (closed and bounded) subset of Euclidean space. Then, for any continuous target function $f \in C(K)$ and any arbitrary error tolerance $\epsilon > 0$, there exists an integer $M \in \mathbb{N}$, input weight vectors $\mathbf{w}_j \in \R^D$, biases $b_j \in \R$, and output weights $\alpha_j \in \R$ such that the single-hidden-layer network:
\begin{equation}
F(\mathbf{x}) = \sum_{j=1}^M \alpha_j \sigma(\mathbf{w}_j^\top \mathbf{x} + b_j),
\label{eq:universal_approximator_form}
\end{equation}
satisfies uniform approximation over all $\mathbf{x} \in K$:
\begin{equation}
\sup_{\mathbf{x} \in K} |F(\mathbf{x}) - f(\mathbf{x})| < \epsilon.
\label{eq:universal_approximation_bound}
\end{equation}
\end{theorem}

\paragraph{Constructive geometric proof via localized bump functions.}
To understand why Theorem~\ref{thm:uat} holds, consider the constructive geometric proof intuition formulated by Cybenko and developed in modern treatments by Bishop~\cite{bishop2024deep} and Prince~\cite{prince2023understanding}.

In one dimension ($D=1$), consider two sigmoid units with identical steep slope $w \gg 0$ but offset biases $b_1 = -w a$ and $b_2 = -w b$ (with $a < b$):
\begin{equation}
\operatorname{bump}(x) \triangleq \sigma(w(x - a)) - \sigma(w(x - b)).
\label{eq:sigmoid_bump}
\end{equation}
As $w \to \infty$, the sigmoid $\sigma(w(x - a))$ converges to a Heaviside step jumping from $0$ to $1$ at $x = a$, while $\sigma(w(x - b))$ jumps from $0$ to $1$ at $x = b$. Subtracting them produces a localized \textbf{rectangular pulse (bump)} that evaluates to approximately $1$ on the interval $[a, b]$ and decays to $0$ everywhere outside.

\begin{figure}[H]
\centering
\begin{tikzpicture}
  \begin{axis}[
    width=13.5cm, height=6.2cm,
    axis lines=center,
    xlabel={$x$},
    ylabel={$y$},
    xmin=-0.5, xmax=6.5,
    ymin=-0.2, ymax=2.2,
    legend pos=north west,
    legend cell align=left,
    legend style={font=\scriptsize, fill=white, fill opacity=0.9},
    tick label style={font=\scriptsize},
    label style={font=\small\sffamily},
    grid=both, grid style={dotted, gray!30}
  ]
    % Target continuous non-linear curve
    \addplot[domain=0:6, samples=150, ultra thick, politoNavy] 
      {0.8 + 0.6*sin(deg(x*1.2)) + 0.3*cos(deg(x*2.1))};
    \addlegendentry{Target Function $f(x) \in C(K)$}

    % Localized sigmoid bumps approximating the curve
    % Bump 1: x in [0.2, 1.4], height 1.3
    \addplot[domain=0:6, samples=200, thick, deepblue, dashed] 
      {1.3 * (1/(1 + exp(-20*(x - 0.2))) - 1/(1 + exp(-20*(x - 1.4))))};
    
    % Bump 2: x in [1.4, 2.8], height 1.65
    \addplot[domain=0:6, samples=200, thick, intuitiongreen, dashed] 
      {1.65 * (1/(1 + exp(-20*(x - 1.4))) - 1/(1 + exp(-20*(x - 2.8))))};

    % Bump 3: x in [2.8, 4.2], height 0.55
    \addplot[domain=0:6, samples=200, thick, deepdiveTeal, dashed] 
      {0.55 * (1/(1 + exp(-20*(x - 2.8))) - 1/(1 + exp(-20*(x - 4.2))))};

    % Bump 4: x in [4.2, 5.8], height 0.95
    \addplot[domain=0:6, samples=200, thick, orange!80!black, dashed] 
      {0.95 * (1/(1 + exp(-20*(x - 4.2))) - 1/(1 + exp(-20*(x - 5.8))))};

    % Total sum of bumps F(x)
    \addplot[domain=0:6, samples=300, line width=1.4pt, examred] 
      {1.3 * (1/(1 + exp(-20*(x - 0.2))) - 1/(1 + exp(-20*(x - 1.4)))) +
       1.65 * (1/(1 + exp(-20*(x - 1.4))) - 1/(1 + exp(-20*(x - 2.8)))) +
       0.55 * (1/(1 + exp(-20*(x - 2.8))) - 1/(1 + exp(-20*(x - 4.2)))) +
       0.95 * (1/(1 + exp(-20*(x - 4.2))) - 1/(1 + exp(-20*(x - 5.8))))};
    \addlegendentry{Network Approximation $F(x) = \sum \alpha_j \operatorname{bump}_j(x)$}

    % Annotations
    \node[font=\tiny\bfseries, text=deepblue] at (axis cs: 0.8, 1.45) {$\alpha_1 \operatorname{bump}_1$};
    \node[font=\tiny\bfseries, text=intuitiongreen] at (axis cs: 2.1, 1.8) {$\alpha_2 \operatorname{bump}_2$};
    \node[font=\tiny\bfseries, text=deepdiveTeal] at (axis cs: 3.5, 0.7) {$\alpha_3 \operatorname{bump}_3$};
    \node[font=\tiny\bfseries, text=orange!80!black] at (axis cs: 5.0, 1.1) {$\alpha_4 \operatorname{bump}_4$};
  \end{axis}
\end{tikzpicture}
\caption{The constructive geometric mechanism of the Universal Approximation Theorem. Pairs of steep sigmoidal neurons subtract to form localized rectangular pulse functions (bumps). By tiling the compact input domain with $M$ localized bumps scaled by coefficients $\alpha_j \approx f(x_j)$, a single-hidden-layer network constructs a Riemann-like step approximation that converges uniformly to any continuous function $f(x)$ as width $M \to \infty$.}
\label{fig:uat_bump_mechanism}
\end{figure}

Any continuous target function $f(x)$ on a compact interval $[x_{\min}, x_{\max}]$ is uniformly continuous. Given tolerance $\epsilon > 0$, we can partition the interval into $M$ sub-intervals of width $\Delta x = \epsilon / L_f$ (where $L_f$ is the Lipschitz constant of $f$). Associating a bump function with each interval $[x_{j-1}, x_j]$ and weighting it by the target function's midpoint value $\alpha_j = f(x_j)$ yields a Riemann step approximation $F(x) = \sum_{j=1}^M \alpha_j \operatorname{bump}_j(x)$. As depicted in Figure~\ref{fig:uat_bump_mechanism}, the approximation error satisfies $|F(x) - f(x)| < \epsilon$ everywhere on the domain.

In $D$-dimensional space $\R^D$, this construction generalizes directly: $2D$ sigmoidal hyperplanes intersect to enclose a localized multidimensional hypercube, forming a localized $D$-dimensional bump. Summing these localized bumps enables the network to tile any compact domain in $\R^D$. For Rectified Linear Units ($\operatorname{ReLU}$), each hidden neuron introduces an affine bend (spline knot); summing $M$ such units constructs a continuous piecewise linear (CPWL) manifold whose planar facets triangulate any arbitrary smooth function.

\paragraph{From shallow approximations to deep feedforward networks.}
Although Theorem~\ref{thm:uat} guarantees that a single hidden layer is theoretically sufficient to approximate any continuous function, practical machine learning almost universally deploys \textbf{deep feedforward neural networks} (also termed \textbf{Fully Connected Networks (FCNs)} or \textbf{Multi-Layer Perceptrons (MLPs)}).

Formally, an $L$-layer Fully Connected Network computes a sequence of coupled affine projections and element-wise non-linear activations:
\begin{align}
\mathbf{a}^{(0)} &\equiv \mathbf{x} \in \R^{d_0}, \label{eq:fcn_input} \\
\mathbf{z}^{(l)} &= \mathbf{W}^{(l)} \mathbf{a}^{(l-1)} + \mathbf{b}^{(l)}, \quad \mathbf{W}^{(l)} \in \R^{d_l \times d_{l-1}}, \, \mathbf{b}^{(l)} \in \R^{d_l}, \quad \forall l \in \{1, \dots, L\}, \label{eq:fcn_preactivation} \\
\mathbf{a}^{(l)} &= \sigma\left( \mathbf{z}^{(l)} \right) \in \R^{d_l}, \label{eq:fcn_activation}
\end{align}
where $\mathbf{z}^{(l)}$ is the layer-$l$ pre-activation vector, $\sigma(\cdot)$ is the non-linear activation function, $\mathbf{a}^{(l)}$ denotes the hidden activations, and the final output is $\hat{\mathbf{y}} \equiv \mathbf{a}^{(L)}$. The total parameter count scales with layer dimensions:
\begin{equation}
N_{\text{dense parameters}} = \sum_{l=1}^L \left( d_{l-1} \cdot d_l + d_l \right).
\label{eq:dense_parameter_count}
\end{equation}

\begin{figure}[H]
\centering
\begin{tikzpicture}[scale=0.95, every node/.style={transform shape}]
  % Styles
  \tikzstyle{inputnode}=[circle, draw=deepblue, fill=deepblue!10, thick, inner sep=2.5pt]
  \tikzstyle{hiddennode1}=[circle, draw=intuitiongreen, fill=intuitiongreen!15, thick, inner sep=2.5pt]
  \tikzstyle{hiddennode2}=[circle, draw=politoNavy, fill=politoNavy!15, thick, inner sep=2.5pt]
  \tikzstyle{outputnode}=[circle, draw=examred, fill=examred!20, thick, inner sep=2.5pt]

  % Layer 0: Input (d0 = 4)
  \foreach \i in {1,...,4} {
    \node[inputnode] (L0_\i) at (0, 2.5 - \i*1.0) {$x_\i$};
  }
  \node[font=\footnotesize\bfseries\color{deepblue}] at (0, 2.2) {Layer $\mathbf{a}^{(0)} \in \R^4$};

  % Layer 1: Hidden 1 (d1 = 5)
  \foreach \j in {1,...,5} {
    \node[hiddennode1] (L1_\j) at (3.0, 3.0 - \j*1.0) {$h^{(1)}_\j$};
  }
  \node[font=\footnotesize\bfseries\color{intuitiongreen}] at (3.0, 2.6) {Layer $\mathbf{a}^{(1)} \in \R^5$};

  % Layer 2: Hidden 2 (d2 = 5)
  \foreach \k in {1,...,5} {
    \node[hiddennode2] (L2_\k) at (6.0, 3.0 - \k*1.0) {$h^{(2)}_\k$};
  }
  \node[font=\footnotesize\bfseries\color{politoNavy}] at (6.0, 2.6) {Layer $\mathbf{a}^{(2)} \in \R^5$};

  % Layer 3: Output (d3 = 3)
  \foreach \m in {1,...,3} {
    \node[outputnode] (L3_\m) at (9.0, 2.0 - \m*1.0) {$\hat{y}_\m$};
  }
  \node[font=\footnotesize\bfseries\color{examred}] at (9.0, 1.6) {Layer $\mathbf{a}^{(3)} \in \R^3$};

  % Connect L0 to L1
  \foreach \i in {1,...,4} {
    \foreach \j in {1,...,5} {
      \draw[-{Stealth}, gray!40, thin] (L0_\i) -- (L1_\j);
    }
  }

  % Connect L1 to L2
  \foreach \j in {1,...,5} {
    \foreach \k in {1,...,5} {
      \draw[-{Stealth}, gray!40, thin] (L1_\j) -- (L2_\k);
    }
  }

  % Connect L2 to L3
  \foreach \k in {1,...,5} {
    \foreach \m in {1,...,3} {
      \draw[-{Stealth}, gray!40, thin] (L2_\k) -- (L3_\m);
    }
  }

  % Annotations for weight matrices
  \node[font=\scriptsize\bfseries\color{gray!80}] at (1.5, -2.5) {$\mathbf{W}^{(1)} \in \R^{5 \times 4}$};
  \node[font=\scriptsize\bfseries\color{gray!80}] at (4.5, -2.5) {$\mathbf{W}^{(2)} \in \R^{5 \times 5}$};
  \node[font=\scriptsize\bfseries\color{gray!80}] at (7.5, -2.5) {$\mathbf{W}^{(3)} \in \R^{3 \times 5}$};
\end{tikzpicture}
\caption{The architectural topology of a deep Fully Connected Network (FCN). Each layer $l$ performs an affine transformation followed by an element-wise non-linear activation: $\mathbf{a}^{(l)} = \sigma(\mathbf{W}^{(l)} \mathbf{a}^{(l-1)} + \mathbf{b}^{(l)})$. In this dense topology, the bipartite graph between adjacent layers is complete: every neuron receives input from all neurons of the preceding layer.}
\label{fig:fcn_topology}
\end{figure}

Why do we transition from wide, single-hidden-layer networks to deep, multi-layer architectures? The fundamental theoretical justification lies in the \textbf{exponential efficiency of depth}:
\begin{enumerate}[noitemsep,topsep=2pt,parsep=0pt,partopsep=0pt]
  \item \textbf{Exponential Linear Regions}: As rigorously proved by Montufar et al. (2014)~\cite{montufar2014number}, a shallow network with $W$ ReLU neurons can partition the input space $\R^D$ into at most $\mathcal{O}(W^D)$ linear regions. In contrast, a deep network with $L$ layers and width $W$ partitions the space into:
  \begin{equation}
  \mathcal{O}\left( \left(\frac{W}{D}\right)^{(L-1)D} W^D \right)
  \label{eq:deep_linear_regions}
  \end{equation}
  linear regions. Depth enables an \textbf{exponential explosion in expressive capacity} for the exact same number of parameters, as subsequent layers repeatedly fold and recombine the piecewise linear boundaries created by earlier layers.
  \item \textbf{Compositional Feature Reuse}: Natural signals (vision, speech, natural language) are governed by hierarchical, compositional physics: pixels combine into edges, edges into motifs and textures, textures into object parts, and parts into complete scenes. Deep architectures embed this hierarchical compositional inductive bias directly into their computational structure: intermediate layers compute reusable features at increasing levels of abstraction, achieving representations with polynomial parameter efficiency $\mathcal{O}(D)$ that would require exponential width $\mathcal{O}(2^D)$ in a shallow network.
\end{enumerate}

\paragraph{Critical drawbacks and limitations of the Universal Approximation Theorem.}
While the Universal Approximation Theorem provides a comforting mathematical guarantee of expressiveness, its practical utility in real-world machine learning is constrained by four fundamental limitations:
\begin{enumerate}[noitemsep,topsep=2pt,parsep=0pt,partopsep=0pt]
  \item \textbf{Existential Guarantee vs. Algorithmic Learnability}: The theorem is purely non-constructive and existential. It proves that there \textit{exists} some choice of width $M$ and parameters $(\mathbf{W}, \mathbf{b}, \boldsymbol{\alpha})$ that achieves an approximation error below $\epsilon$. Crucially, it provides \textbf{zero guarantee that gradient-based optimization algorithms can ever discover these parameters}. The empirical loss landscape of an ultra-wide shallow network is characterized by pathological saddle points, plateaus, and bad local minima that prevent standard gradient descent from converging to the theoretical solution.
  \item \textbf{The Curse of Dimensionality and Exponential Width}: The constructive proof partitions each input dimension into $\mathcal{O}(1/\epsilon)$ intervals. To tile a $D$-dimensional compact unit hypercube $[0, 1]^D$, the required number of localized bump units scales exponentially with dimension:
  \begin{equation}
  M \sim \mathcal{O}\left( \left(\frac{1}{\epsilon}\right)^D \right).
  \label{eq:exponential_width_curse}
  \end{equation}
  For high-dimensional inputs typical of modern tasks (e.g., $D = 3 \times 224 \times 224 \approx 1.5 \times 10^5$ for image classification), $M$ astronomical exceeds the number of elementary particles in the observable universe. Shallow universal approximation is therefore computationally and physically impossible in high dimensions.
  \item \textbf{Zero Generalization Guarantees and Lookup-Table Memorization}: The theorem guarantees uniform approximation only on the compact training domain $K$; it offers \textbf{no generalization bound} on unseen test distributions. An ultra-wide shallow network with $M \to \infty$ units can act as a brute-force lookup table or memorization spline, dedicating individual localized bumps to interpolate every training point exactly. Such a model achieves zero training error but learns zero underlying structure, leading to catastrophic overfitting and failure to generalize.
  \item \textbf{Exponential Sample Complexity}: Accurately configuring the weights $\alpha_j \approx f(\mathbf{x}_j)$ across $M \sim \mathcal{O}((1/\epsilon)^D)$ localized bump functions requires observing training samples in virtually every sub-cube of the domain. In high dimensions, the volume of space is vast and finite training sets leave the input space virtually empty (the empty space phenomenon). Without the strong compositional inductive biases enforced by deep hierarchical architectures, a shallow network requires an exponentially large sample size $N \sim \mathcal{O}((1/\epsilon)^D)$ to avoid arbitrary interpolation errors in unobserved regions.
\end{enumerate}

Yet, while deep feedforward networks possess extraordinary theoretical expressive power, an urgent operational crisis emerges: \textit{how do we compute the parameter updates for a deep network containing millions of entangled weights?} In the next section, we discover why classical algebraic methods fail and how the computational graph resolves this crisis.

\section{Computational graphs and backpropagation}
\label{sec:computational_graphs_autodiff}
\label{sec:backpropagation_derivation}

Having established that deep feedforward networks possess the expressive capacity to represent complex non-linear manifolds, we confront the central operational challenge of machine learning: \textbf{parameter optimization}. In classical Ordinary Least Squares (Chapter~3), parameters admit an exact, globally optimal analytical solution via the Normal Equations $\mathbf{w}^* = (\mathbf{X}^\top \mathbf{X})^{-1}\mathbf{X}^\top \mathbf{y}$. In deep neural networks, however, parameters are deeply entangled across multi-layer compositions of non-linear activations. Optimizing these models requires computing the gradient of a scalar loss function with respect to millions or billions of parameters efficiently and accurately.

\subsection{Motivation for computational graphs}
\label{sec:algebraic_impasse}
\label{sec:numerical_differentiation_failure}

\paragraph{The algebraic impasse and non-convexity.}
A natural instinct is to ask whether an analytical closed-form solution exists for deep neural networks, analogous to the Normal Equations of linear regression. Consider an $L$-layer neural network hypothesis evaluating an input $\mathbf{x} \in \R^{d_0}$:
\begin{equation}
f(\mathbf{x}; \boldsymbol{\theta}) = \mathbf{W}^{(L)} \sigma\left( \mathbf{W}^{(L-1)} \sigma\left( \dots \sigma\left( \mathbf{W}^{(1)} \mathbf{x} + \mathbf{b}^{(1)} \right) \dots \right) + \mathbf{b}^{(L-1)} \right) + \mathbf{b}^{(L)},
\label{eq:deep_composition_hypothesis}
\end{equation}
evaluated across training dataset $\mathcal{D} = \{(\mathbf{x}_i, y_i)\}_{i=1}^N$ under the empirical risk objective:
\begin{equation}
\hat{\mathcal{R}}_N(\boldsymbol{\theta}) = \frac{1}{N} \sum_{i=1}^N \mathcal{L}\left( f(\mathbf{x}_i; \boldsymbol{\theta}), \, y_i \right).
\end{equation}
To obtain an exact analytical minimizer, one must equate the gradient with respect to all parameter matrices to zero: $\nabla_{\boldsymbol{\theta}} \hat{\mathcal{R}}_N(\boldsymbol{\theta}) = \mathbf{0}$. This stationary condition defines a massive system of coupled, transcendental non-linear algebraic equations. Because early parameters $\mathbf{W}^{(1)}$ are deeply nested within cascading non-linear activation functions $\sigma(\cdot)$, they cannot be decoupled, factorized, or isolated through matrix inversion or linear algebraic operations. Furthermore, as established in Chapter~3, the loss surface of deep networks is non-convex, characterized by uncountably many saddle points, local minima, and permutation symmetries. Analytical closed-form solutions are mathematically impossible, necessitating iterative numerical optimization via gradient descent.

\paragraph{The failure of numerical differentiation.}
To execute gradient descent, we must compute the exact gradient vector $\nabla_{\boldsymbol{\theta}} \mathcal{L}$ for networks containing millions to billions of parameters ($W \sim 10^7$--$10^{11}$). The simplest numerical approach for computing derivatives is \textbf{finite differences} (numerical perturbation):
\begin{equation}
\frac{\partial \mathcal{L}}{\partial \theta_j} \approx \frac{\mathcal{L}(\boldsymbol{\theta} + \epsilon \mathbf{e}_j) - \mathcal{L}(\boldsymbol{\theta})}{\epsilon},
\label{eq:finite_differences}
\end{equation}
where $\mathbf{e}_j$ is the $j$-th standard basis vector and $\epsilon \approx 10^{-5}$ is an infinitesimal perturbation. Evaluating Equation~\eqref{eq:finite_differences} across all $W$ parameters requires evaluating $W + 1$ complete forward passes through the network. In a standard deep network where $W = 10^8$ (such as ResNet or AlexNet), computing a single parameter update step would require $100{,}000{,}001$ forward network evaluations. If a single forward pass takes $10$ milliseconds, evaluating one gradient step would require over \textbf{11 days of continuous GPU computation}, yielding an intractable computational complexity of $\mathcal{O}(W^2)$. Finite-difference numerical differentiation is therefore completely unusable for deep learning optimization.

\subsection{DAGs and reverse-mode differentiation}
\label{sec:dag_computational_graphs}
\label{sec:autodiff_modes}

\paragraph{The general principle of computational DAGs.}
Modern deep learning frameworks (such as PyTorch, JAX, and TensorFlow) resolve the differentiation impasse by reformulating network evaluation as a \textbf{Directed Acyclic Graph (DAG)}. In a computational DAG:
\begin{itemize}[noitemsep,topsep=2pt,parsep=0pt,partopsep=0pt]
  \item \textbf{Nodes} represent mathematical variables (tensors: scalars, vectors, matrices, or higher-order arrays).
  \item \textbf{Directed Edges} represent primitive, differentiable elementary operations (such as addition $+$, scalar/matrix multiplication $\times$, element-wise activation $\sigma$, and logarithms $\ln$).
\end{itemize}
Because the graph is acyclic, its nodes admit a \textbf{topological sorting} that defines a valid execution order. Evaluating the network on the DAG naturally decomposes into two phases: (1) a \textbf{Forward Pass}, where inputs flow along directed edges according to topological sort order, computing and caching intermediate activations at each node; and (2) a \textbf{Backward Pass}, where the chain rule of differential calculus propagates derivative signals in reverse topological order, modulating incoming error signals by local Jacobian matrices.

\begin{figure}[H]
\centering
\begin{tikzpicture}[scale=1.05]
  % Nodes
  \node[circle, draw=deepblue, fill=deepblue!15, thick, inner sep=3pt] (x) at (0, 1.2) {$x$};
  \node[circle, draw=deepblue, fill=deepblue!15, thick, inner sep=3pt] (w) at (0, -0.2) {$w$};
  \node[circle, draw=politoNavy, fill=politoNavy!20, thick, inner sep=3pt] (mul) at (2.0, 0.5) {$\times$};
  \node[circle, draw=deepblue, fill=deepblue!15, thick, inner sep=3pt] (b) at (2.0, -1.0) {$b$};
  \node[circle, draw=politoNavy, fill=politoNavy!20, thick, inner sep=3pt] (add) at (4.0, 0.0) {$+$};
  \node[circle, draw=intuitiongreen, fill=intuitiongreen!20, thick, inner sep=3pt] (sig) at (6.0, 0.0) {$\sigma$};
  \node[circle, draw=deepblue, fill=deepblue!15, thick, inner sep=3pt] (y) at (6.0, -1.2) {$y$};
  \node[circle, draw=examred, fill=examred!25, thick, inner sep=3pt] (loss) at (8.0, -0.6) {$\mathcal{L}$};

  % Forward Edges (Blue)
  \draw[-{Stealth}, very thick, deepblue] (x) -- node[above, font=\scriptsize] {$x$} (mul);
  \draw[-{Stealth}, very thick, deepblue] (w) -- node[below, font=\scriptsize] {$w$} (mul);
  \draw[-{Stealth}, very thick, deepblue] (mul) -- node[above, font=\scriptsize] {$wx$} (add);
  \draw[-{Stealth}, very thick, deepblue] (b) -- node[below, font=\scriptsize] {$b$} (add);
  \draw[-{Stealth}, very thick, deepblue] (add) -- node[above, font=\scriptsize] {$z$} (sig);
  \draw[-{Stealth}, very thick, deepblue] (sig) -- node[above, font=\scriptsize] {$\hat{y}$} (loss);
  \draw[-{Stealth}, very thick, deepblue] (y) -- node[below, font=\scriptsize] {$y$} (loss);

  % Backward Adjoint Flow (Red Dashed)
  \draw[-{Stealth}, thick, examred, dashed] (loss) to[bend left=30] node[above, font=\tiny\bfseries] {$\frac{\partial \mathcal{L}}{\partial \hat{y}}$} (sig);
  \draw[-{Stealth}, thick, examred, dashed] (sig) to[bend left=30] node[above, font=\tiny\bfseries] {$\frac{\partial \mathcal{L}}{\partial z}$} (add);
  \draw[-{Stealth}, thick, examred, dashed] (add) to[bend left=30] node[above, font=\tiny\bfseries] {$\frac{\partial \mathcal{L}}{\partial (wx)}$} (mul);
  \draw[-{Stealth}, thick, examred, dashed] (add) to[bend left=30] node[below, font=\tiny\bfseries] {$\frac{\partial \mathcal{L}}{\partial b}$} (b);
  \draw[-{Stealth}, thick, examred, dashed] (mul) to[bend left=30] node[below, font=\tiny\bfseries] {$\frac{\partial \mathcal{L}}{\partial w}$} (w);

  % Legend
  \node[draw=gray!60, fill=white, rounded corners, font=\scriptsize, align=left] at (4.0, 2.0) {
    \textcolor{deepblue}{\textbf{Forward Pass}}: Activations evaluated and cached\\
    \textcolor{examred}{\textbf{Backward Pass}}: Adjoint sensitivities propagated via chain rule
  };
\end{tikzpicture}
\caption{A Directed Acyclic Graph (DAG) representation of a computational graph for an artificial neuron with scalar loss $\mathcal{L}(\sigma(wx + b), y)$. Forward evaluation (solid blue lines) evaluates and caches intermediate variables. Reverse-mode automatic differentiation (dashed red lines) propagates adjoint sensitivities backwards from the scalar loss.}
\label{fig:computational_graph}
\end{figure}

\paragraph{Reverse-mode automatic differentiation.}
Let $f: \R^n \to \R^m$ be an arbitrary composite differentiable function executed on a DAG. Automatic differentiation computes exact derivatives at machine precision by decomposing $f$ into elementary operations and systematically applying the chain rule. It operates in one of two fundamental orientations:
\begin{enumerate}[noitemsep,topsep=2pt,parsep=0pt,partopsep=0pt]
  \item \textbf{Forward Tangent Mode}: Tracks directional derivatives (tangents) alongside intermediate variables during the forward pass: $\dot{v}_i \triangleq \sum_{j=1}^n \frac{\partial v_i}{\partial x_j} \dot{x}_j$. To compute the full Jacobian $\mathbf{J} \in \R^{m \times n}$, forward-mode autodiff requires seeding the input tangents with standard basis vectors $\dot{\mathbf{x}} = \mathbf{e}_k$ across $n$ independent forward sweeps. Forward mode is computationally optimal when the number of inputs is much smaller than the number of outputs ($n \ll m$).
  \item \textbf{Reverse Adjoint Mode}: Executes in two sweeps. First, the forward sweep evaluates and caches all intermediate node values $v_i$. Second, beginning at the terminal output nodes, it propagates adjoint sensitivity variables backward:
  \begin{equation}
  \bar{v}_i \triangleq \frac{\partial \mathcal{L}}{\partial v_i} = \sum_{j \in \operatorname{children}(i)} \bar{v}_j \frac{\partial v_j}{\partial v_i}.
  \label{eq:adjoint_definition}
  \end{equation}
  To evaluate the complete Jacobian, reverse mode seeds the output adjoints $\bar{\mathbf{y}} = \mathbf{e}_k$ and sweeps backward $m$ times.
\end{enumerate}

In deep learning, the training objective is a \textbf{scalar loss} ($\mathcal{L} \in \R$, so output dimension $m = 1$), while the parameter space is massive ($n = W \sim 10^8$). Because $m = 1$, reverse-mode automatic differentiation evaluates the \textbf{entire parameter gradient vector $\nabla_{\boldsymbol{\theta}} \mathcal{L} \in \R^W$ in exactly ONE single backward pass}! It reduces computational complexity from $\mathcal{O}(W^2)$ to $\mathcal{O}(W)$, slashing gradient evaluation time by eight orders of magnitude.

\begin{table}[H]
\centering
\small
\setlength{\tabcolsep}{6pt}
\caption{Comparison of differentiation mechanisms across algorithmic paradigms.}
\label{tab:autodiff_comparison}
\vspace{4pt}
\begin{tabularx}{\textwidth}{l c c c >{\raggedright\arraybackslash}X}
\toprule
\textbf{Method} & \textbf{Exactness} & \textbf{Complexity} & \textbf{Memory} & \textbf{Suitability for Deep Learning} \\
\midrule
\textbf{Numerical (Finite Diff.)} & Approximate ($\mathcal{O}(\epsilon)$) & $\mathcal{O}(W^2)$ & $\mathcal{O}(1)$ & Unusable: $10^8$ forward passes per step. \\
\textbf{Symbolic (Algebraic)} & Exact & Exponential & Memory blowup & Unusable: expression swell on deep DAGs. \\
\textbf{Forward-Mode Autodiff} & Exact & $\mathcal{O}(n \cdot W)$ & $\mathcal{O}(1)$ & Sub-optimal: requires $W$ passes for scalar loss. \\
\textbf{Reverse-Mode Autodiff} & Exact & $\mathcal{O}(m \cdot W)$ & $\mathcal{O}(W)$ & \textbf{Ideal}: exactly \textbf{ONE backward pass} for scalar loss ($m=1$). \\
\bottomrule
\end{tabularx}
\end{table}

\subsection{The backpropagation algorithm}
\label{sec:four_equations_backprop}
\label{sec:vectorized_minibatch_backprop}
\label{sec:memory_footprint_caching}
\label{sec:backprop_numerical_walkthrough}

\paragraph{The general idea behind backpropagation.}
The Backpropagation algorithm is the concrete, layer-wise realization of reverse-mode automatic differentiation applied to neural network architectures. Formulated in optimal control by Paul Werbos (1974)~\cite{werbos1974beyond} and popularized in connectionist machine learning by David Rumelhart, Geoffrey Hinton, and Ronald Williams (1986)~\cite{rumelhart1986learning}, backpropagation represents the foundational algorithmic engine of modern deep learning.

The central intuition of backpropagation is solving the \textbf{credit assignment problem}: determining how much each individual weight and bias in an arbitrary layer contributed to the final prediction error. Instead of computing derivatives independently from scratch for every parameter, backpropagation transforms credit assignment into a localized, recursive message-passing scheme. Information flows in two directions: forward signals compute intermediate representations, while backward signals (error adjoints) propagate from the loss back toward the inputs. Each layer receives an error message from the layer above, uses its cached forward activations to compute its own parameter gradients, and multiplies the error by its transposed weight matrix to pass an error message to the layer below.

\paragraph{The algorithmic cycle of backpropagation.}
The execution of backpropagation across an $L$-layer network proceeds through five sequential steps:
\begin{enumerate}[noitemsep,topsep=2pt,parsep=0pt,partopsep=0pt]
  \item \textbf{Forward Propagation}: Feed input $\mathbf{x} \equiv \mathbf{a}^{(0)}$ through all layers $l = 1, \dots, L$. At each layer, compute and cache pre-activations $\mathbf{z}^{(l)} = \mathbf{W}^{(l)} \mathbf{a}^{(l-1)} + \mathbf{b}^{(l)}$ and post-activations $\mathbf{a}^{(l)} = \sigma(\mathbf{z}^{(l)})$.
  \item \textbf{Loss and Output Error Evaluation}: Evaluate the scalar loss $\mathcal{L}(\mathbf{a}^{(L)}, \mathbf{y})$ and compute the output layer error adjoint $\boldsymbol{\delta}^{(L)} = \frac{\partial \mathcal{L}}{\partial \mathbf{z}^{(L)}}$.
  \item \textbf{Backward Error Propagation}: Recursively backpropagate the error adjoints $\boldsymbol{\delta}^{(l)}$ from layer $l+1$ back to layer $l$ for $l = L-1, L-2, \dots, 1$ via transposed matrix multiplication and element-wise activation differentiation.
  \item \textbf{Parameter Gradient Accumulation}: Compute partial derivatives of the loss with respect to all parameter matrices $\frac{\partial \mathcal{L}}{\partial \mathbf{W}^{(l)}}$ and bias vectors $\frac{\partial \mathcal{L}}{\partial \mathbf{b}^{(l)}}$ using the error adjoints $\boldsymbol{\delta}^{(l)}$ and stored activations $\mathbf{a}^{(l-1)}$.
  \item \textbf{Parameter Update}: Update all network parameters via gradient descent: $\mathbf{W}^{(l)} \leftarrow \mathbf{W}^{(l)} - \eta \frac{\partial \mathcal{L}}{\partial \mathbf{W}^{(l)}}$ and $\mathbf{b}^{(l)} \leftarrow \mathbf{b}^{(l)} - \eta \frac{\partial \mathcal{L}}{\partial \mathbf{b}^{(l)}}$.
\end{enumerate}

\paragraph{Step-by-step mathematical derivation of the four fundamental equations.}
We define the \textbf{layer error adjoint vector} $\boldsymbol{\delta}^{(l)}$ as the partial derivative of the scalar loss with respect to the pre-activation vector $\mathbf{z}^{(l)}$:
\begin{equation}
\boldsymbol{\delta}^{(l)} \triangleq \frac{\partial \mathcal{L}}{\partial \mathbf{z}^{(l)}} \in \R^{d_l}.
\label{eq:delta_definition_ch4}
\end{equation}
The vector $\boldsymbol{\delta}^{(l)}$ quantifies the sensitivity of the final loss to perturbations in layer $l$'s pre-activations. The algorithm is governed by four fundamental equations:

\begin{enumerate}[noitemsep,topsep=2pt,parsep=0pt,partopsep=0pt]
  \item \textbf{Equation BP1: Output Layer Error Adjoint $\boldsymbol{\delta}^{(L)}$.}
  Applying the multivariate chain rule to the scalar loss $\mathcal{L}$ with respect to coordinate $z_k^{(L)}$:
  \begin{equation}
  \delta_k^{(L)} = \frac{\partial \mathcal{L}}{\partial z_k^{(L)}} = \sum_{j=1}^{d_L} \frac{\partial \mathcal{L}}{\partial a_j^{(L)}} \frac{\partial a_j^{(L)}}{\partial z_k^{(L)}}.
  \end{equation}
  Because the activation function $\sigma(\cdot)$ is applied element-wise, $\frac{\partial a_j^{(L)}}{\partial z_k^{(L)}} = 0$ for all $j \ne k$. The summation collapses to a single term:
  \begin{equation}
  \delta_k^{(L)} = \frac{\partial \mathcal{L}}{\partial a_k^{(L)}} \sigma'\left( z_k^{(L)} \right) \implies \boldsymbol{\delta}^{(L)} = \nabla_{\mathbf{a}^{(L)}} \mathcal{L} \odot \sigma'\left( \mathbf{z}^{(L)} \right),
  \label{eq:bp1_output_error}
  \end{equation}
  where $\odot$ denotes the Hadamard (element-wise) product. For standard Cross-Entropy loss paired with Softmax or Sigmoid activations, this simplifies directly to the residual prediction error: $\boldsymbol{\delta}^{(L)} = \hat{\mathbf{y}} - \mathbf{y}$.

  \item \textbf{Equation BP2: Hidden Layer Error Recurrence $\boldsymbol{\delta}^{(l)}$.}
  Suppose the error adjoint vector $\boldsymbol{\delta}^{(l+1)} \in \R^{d_{l+1}}$ at layer $l+1$ has already been computed. By the multivariate chain rule:
  \begin{equation}
  \delta_j^{(l)} = \frac{\partial \mathcal{L}}{\partial z_j^{(l)}} = \sum_{k=1}^{d_{l+1}} \frac{\partial \mathcal{L}}{\partial z_k^{(l+1)}} \frac{\partial z_k^{(l+1)}}{\partial z_j^{(l)}} = \sum_{k=1}^{d_{l+1}} \delta_k^{(l+1)} \frac{\partial z_k^{(l+1)}}{\partial z_j^{(l)}}.
  \label{eq:delta_chain_component_ch4}
  \end{equation}
  Since $z_k^{(l+1)} = \sum_{m=1}^{d_l} W_{km}^{(l+1)} \sigma(z_m^{(l)}) + b_k^{(l+1)}$, differentiating with respect to $z_j^{(l)}$ yields $\frac{\partial z_k^{(l+1)}}{\partial z_j^{(l)}} = W_{kj}^{(l+1)} \sigma'(z_j^{(l)})$. Substituting into Equation~\eqref{eq:delta_chain_component_ch4}:
  \begin{equation}
  \delta_j^{(l)} = \sum_{k=1}^{d_{l+1}} \delta_k^{(l+1)} W_{kj}^{(l+1)} \sigma'\left( z_j^{(l)} \right) = \left( \sum_{k=1}^{d_{l+1}} W_{kj}^{(l+1)} \delta_k^{(l+1)} \right) \sigma'\left( z_j^{(l)} \right).
  \end{equation}
  In vector notation, this produces the core backpropagation recurrence:
  \begin{equation}
  \boldsymbol{\delta}^{(l)} = \left( (\mathbf{W}^{(l+1)})^\top \boldsymbol{\delta}^{(l+1)} \right) \odot \sigma'\left( \mathbf{z}^{(l)} \right).
  \label{eq:bp2_error_recurrence}
  \end{equation}
  The error adjoint propagates backward by multiplying with the transposed weight matrix $(\mathbf{W}^{(l+1)})^\top$, followed by element-wise modulation by the local activation derivative $\sigma'(\mathbf{z}^{(l)})$.

  \item \textbf{Equation BP3: Weight Matrix Parameter Gradients.}
  The gradient of the loss with respect to individual weight entries $W_{jk}^{(l)}$ is:
  \begin{equation}
  \frac{\partial \mathcal{L}}{\partial W_{jk}^{(l)}} = \frac{\partial \mathcal{L}}{\partial z_j^{(l)}} \frac{\partial z_j^{(l)}}{\partial W_{jk}^{(l)}} = \delta_j^{(l)} a_k^{(l-1)}.
  \end{equation}
  In outer-product matrix notation:
  \begin{equation}
  \frac{\partial \mathcal{L}}{\partial \mathbf{W}^{(l)}} = \boldsymbol{\delta}^{(l)} (\mathbf{a}^{(l-1)})^\top \in \R^{d_l \times d_{l-1}}.
  \label{eq:grad_weight_matrix}
  \end{equation}

  \item \textbf{Equation BP4: Bias Vector Parameter Gradients.}
  Similarly, for the scalar bias parameter $b_j^{(l)}$:
  \begin{equation}
  \frac{\partial \mathcal{L}}{\partial b_j^{(l)}} = \frac{\partial \mathcal{L}}{\partial z_j^{(l)}} \frac{\partial z_j^{(l)}}{\partial b_j^{(l)}} = \delta_j^{(l)} \cdot 1 \implies \frac{\partial \mathcal{L}}{\partial \mathbf{b}^{(l)}} = \boldsymbol{\delta}^{(l)} \in \R^{d_l}.
  \label{eq:bp4_bias_gradient}
  \end{equation}
\end{enumerate}

In practical machine learning, computations are vectorized over mini-batches of size $B$. Stacking $B$ training observations as rows in matrix $\mathbf{X} \in \R^{B \times d_0}$, the batched forward pass computes:
\begin{equation}
\mathbf{Z}^{(l)} = \mathbf{A}^{(l-1)} (\mathbf{W}^{(l)})^\top + \mathbf{1}_B (\mathbf{b}^{(l)})^\top, \quad \mathbf{A}^{(l)} = \sigma\left( \mathbf{Z}^{(l)} \right),
\label{eq:batched_forward_equations}
\end{equation}
where $\mathbf{1}_B = [1, \dots, 1]^\top \in \R^B$ broadcasts the bias row vector across all samples. Defining the batch error adjoint matrix $\boldsymbol{\Delta}^{(l)} \triangleq \nabla_{\mathbf{Z}^{(l)}} \mathcal{L} \in \R^{B \times d_l}$, the vectorized batched backpropagation equations evaluate to:
\begin{align}
\boldsymbol{\Delta}^{(L)} &= \nabla_{\mathbf{A}^{(L)}} \mathcal{L} \odot \sigma'\left( \mathbf{Z}^{(L)} \right) \in \R^{B \times d_L}, \label{eq:batched_bp1} \\
\boldsymbol{\Delta}^{(l)} &= \left( \boldsymbol{\Delta}^{(l+1)} \mathbf{W}^{(l+1)} \right) \odot \sigma'\left( \mathbf{Z}^{(l)} \right) \in \R^{B \times d_l}, \label{eq:batched_bp2} \\
\frac{\partial \mathcal{L}}{\partial \mathbf{W}^{(l)}} &= \frac{1}{B} \left( \boldsymbol{\Delta}^{(l)} \right)^\top \mathbf{A}^{(l-1)} \in \R^{d_l \times d_{l-1}}, \label{eq:batched_bp3} \\
\frac{\partial \mathcal{L}}{\partial \mathbf{b}^{(l)}} &= \frac{1}{B} \left( \boldsymbol{\Delta}^{(l)} \right)^\top \mathbf{1}_B \in \R^{d_l}. \label{eq:batched_bp4}
\end{align}

\begin{figure}[H]
\centering
\begin{tikzpicture}[scale=0.95, every node/.style={transform shape}]
  % Styles
  \tikzstyle{databox}=[rectangle, draw=deepblue, fill=deepblue!10, rounded corners=4pt, thick, align=center, font=\footnotesize\sffamily]
  \tikzstyle{weightbox}=[rectangle, draw=politoNavy, fill=politoNavy!15, rounded corners=4pt, thick, align=center, font=\footnotesize\sffamily]
  \tikzstyle{opbox}=[circle, draw=intuitiongreen, fill=intuitiongreen!20, thick, font=\small]
  \tikzstyle{lossbox}=[rectangle, draw=examred, fill=examred!20, rounded corners=4pt, very thick, align=center, font=\footnotesize\bfseries\sffamily]

  % Forward Nodes (Center Line y = 1.8)
  \node[databox] (X) at (0, 1.8) {$\mathbf{X} \equiv \mathbf{A}^{(0)}$\\$\R^{B \times d_0}$};
  \node[databox] (Z1) at (3.0, 1.8) {$\mathbf{Z}^{(1)}$\\$\R^{B \times d_1}$};
  \node[opbox] (sig1) at (4.5, 1.8) {$\sigma$};
  \node[databox] (A1) at (6.0, 1.8) {$\mathbf{A}^{(1)}$\\$\R^{B \times d_1}$};
  \node[databox] (Z2) at (9.0, 1.8) {$\mathbf{Z}^{(2)}$\\$\R^{B \times d_2}$};
  \node[lossbox] (loss) at (11.5, 1.8) {Loss $\mathcal{L}$\\$\R$ (scalar)};

  % Parameter Nodes (Placed Above at y = 3.6)
  \node[weightbox] (W1) at (3.0, 3.6) {$\mathbf{W}^{(1)}$\\$\R^{d_1 \times d_0}$};
  \node[weightbox] (W2) at (9.0, 3.6) {$\mathbf{W}^{(2)}$\\$\R^{d_2 \times d_1}$};

  % Forward Edges (Solid Blue)
  \draw[-{Stealth}, very thick, deepblue] (X) -- (Z1) node[midway, above, font=\tiny] {$\mathbf{X}\mathbf{W}^{(1)\top}$};
  \draw[-{Stealth}, thick, deepblue] (W1) -- (Z1);
  \draw[-{Stealth}, very thick, deepblue] (Z1) -- (sig1);
  \draw[-{Stealth}, very thick, deepblue] (sig1) -- (A1);
  \draw[-{Stealth}, very thick, deepblue] (A1) -- (Z2) node[midway, above, font=\tiny] {$\mathbf{A}^{(1)}\mathbf{W}^{(2)\top}$};
  \draw[-{Stealth}, thick, deepblue] (W2) -- (Z2);
  \draw[-{Stealth}, very thick, deepblue] (Z2) -- (loss);

  % Backward Adjoint Flows (Dashed Red, Bottom Lane y <= 0.6)
  \draw[-{Stealth}, very thick, examred, dashed] (loss.south) to[bend left=30] node[below, font=\tiny\bfseries] {$\boldsymbol{\Delta}^{(2)} \in \R^{B \times d_2}$} (Z2.south);
  \draw[-{Stealth}, very thick, examred, dashed] (Z2.south) to[bend left=30] node[below, font=\tiny\bfseries] {$\boldsymbol{\Delta}^{(2)}\mathbf{W}^{(2)}$} (A1.south);
  \draw[-{Stealth}, very thick, examred, dashed] (A1.south) to[bend left=30] node[below, font=\tiny\bfseries] {$\odot \sigma'(\mathbf{Z}^{(1)})$} (Z1.south);
  \draw[-{Stealth}, very thick, examred, dashed] (Z1.south) to[bend left=30] node[below, font=\tiny\bfseries] {$\boldsymbol{\Delta}^{(1)} \in \R^{B \times d_1}$} (X.south);

  % Parameter Gradient Updates (Upward Red Dashed Arrows to Weights)
  \draw[-{Stealth}, thick, examred, dashed] (Z1.north) to[bend right=25] node[left, font=\tiny\bfseries] {$\frac{1}{B}\boldsymbol{\Delta}^{(1)\top}\mathbf{X}$} (W1.south);
  \draw[-{Stealth}, thick, examred, dashed] (Z2.north) to[bend right=25] node[left, font=\tiny\bfseries] {$\frac{1}{B}\boldsymbol{\Delta}^{(2)\top}\mathbf{A}^{(1)}$} (W2.south);

  % Legend
  \node[draw=gray!50, fill=white, rounded corners, font=\scriptsize, align=left] at (6.0, -0.6) {
    \textcolor{deepblue}{\textbf{Forward Pass}}: Tensors evaluated and activations cached in GPU VRAM\\
    \textcolor{examred}{\textbf{Backward Pass}}: Adjoint error matrices $\boldsymbol{\Delta}^{(l)}$ and parameter gradients propagated via GEMM contractions
  };
\end{tikzpicture}
\caption{Batched tensor computational graph for a multi-layer feedforward network. Forward execution (solid blue) propagates activation matrices $\mathbf{A}^{(l)} \in \mathbb{R}^{B \times d_l}$ through weight matrices. Backward execution (dashed red) propagates adjoint error matrices $\boldsymbol{\Delta}^{(l)} \in \mathbb{R}^{B \times d_l}$, computing parameter gradient matrices via outer-product matrix contractions.}
\label{fig:batched_computational_graph}
\end{figure}

In Equation~\eqref{eq:batched_bp3}, the dimensions contract as $(d_l \times B) \times (B \times d_{l-1}) = (d_l \times d_{l-1})$. The inner batch dimension $B$ contracts completely, accumulating the sample-averaged outer product across all $B$ observations in a single high-performance BLAS GEMM matrix multiplication!

\paragraph{Concrete end-to-end numerical walkthrough.}
To make the algorithm completely transparent and provide a concrete numerical benchmark, consider a minimal two-layer network evaluated on a single observation:
\begin{itemize}[noitemsep,topsep=2pt,parsep=0pt,partopsep=0pt]
  \item Input: $\mathbf{x} = [x_1, x_2]^\top = [0.5, -0.2]^\top \in \R^2$.
  \item Hidden layer: 2 neurons with logistic sigmoid activation $\sigma(z) = \frac{1}{1 + e^{-z}}$, weights $\mathbf{W}^{(1)} = \begin{bmatrix} 0.4 & -0.5 \\ 0.2 & 0.8 \end{bmatrix}$, and bias $\mathbf{b}^{(1)} = \begin{bmatrix} 0.1 \\ -0.1 \end{bmatrix}$.
  \item Output layer: 1 linear neuron $\hat{y} = z^{(2)} = \mathbf{w}^{(2)\top} \mathbf{a}^{(1)} + b^{(2)}$ with $\mathbf{w}^{(2)} = \begin{bmatrix} 0.7 \\ -0.3 \end{bmatrix}$, $b^{(2)} = 0.2$, target $y = 1.0$, and squared error loss $\mathcal{L} = \frac{1}{2}(\hat{y} - y)^2$.
\end{itemize}

\textit{1. Forward Pass Evaluation:}
The hidden pre-activations evaluate to:
\begin{align}
z_1^{(1)} &= (0.4)(0.5) + (-0.5)(-0.2) + 0.1 = 0.20 + 0.10 + 0.1 = 0.40, \nonumber \\
z_2^{(1)} &= (0.2)(0.5) + (0.8)(-0.2) - 0.1 = 0.10 - 0.16 - 0.1 = -0.16.
\end{align}
Applying the logistic sigmoid activation:
\begin{align}
a_1^{(1)} &= \sigma(0.40) = \frac{1}{1 + e^{-0.40}} \approx 0.5987, \nonumber \\
a_2^{(1)} &= \sigma(-0.16) = \frac{1}{1 + e^{0.16}} \approx 0.4601.
\end{align}
Evaluating the output pre-activation and prediction:
\begin{align}
\hat{y} = z^{(2)} &= (0.7)(0.5987) + (-0.3)(0.4601) + 0.2 \approx 0.4191 - 0.1380 + 0.2 = 0.4811.
\end{align}
The forward loss is:
\begin{equation}
\mathcal{L} = \frac{1}{2}(\hat{y} - y)^2 = \frac{1}{2}(0.4811 - 1.0)^2 = \frac{1}{2}(-0.5189)^2 \approx 0.1346.
\end{equation}

\textit{2. Backward Pass Evaluation:}
Step 1: Compute the output layer error adjoint $\delta^{(2)}$ via Equation BP1:
\begin{equation}
\delta^{(2)} \triangleq \frac{\partial \mathcal{L}}{\partial z^{(2)}} = \hat{y} - y = 0.4811 - 1.0 = -0.5189.
\end{equation}
The output parameter gradients follow immediately via Equations BP3 and BP4:
\begin{align}
\frac{\partial \mathcal{L}}{\partial w_1^{(2)}} &= \delta^{(2)} a_1^{(1)} = (-0.5189)(0.5987) \approx -0.3107, \nonumber \\
\frac{\partial \mathcal{L}}{\partial w_2^{(2)}} &= \delta^{(2)} a_2^{(1)} = (-0.5189)(0.4601) \approx -0.2387, \nonumber \\
\frac{\partial \mathcal{L}}{\partial b^{(2)}} &= \delta^{(2)} = -0.5189.
\end{align}

Step 2: Backpropagate error adjoints to the hidden layer pre-activations via Equation BP2:
\begin{equation}
\boldsymbol{\delta}^{(1)} = \left( (\mathbf{w}^{(2)})^\top \delta^{(2)} \right) \odot \sigma'(\mathbf{z}^{(1)}).
\end{equation}
Evaluating the sigmoid derivatives $\sigma'(z) = \sigma(z)(1 - \sigma(z))$:
\begin{align}
\sigma'(z_1^{(1)}) &= (0.5987)(1 - 0.5987) \approx 0.2403, \nonumber \\
\sigma'(z_2^{(1)}) &= (0.4601)(1 - 0.4601) \approx 0.2484.
\end{align}
Multiplying the backpropagated error by the activation derivatives:
\begin{align}
\delta_1^{(1)} &= (w_1^{(2)} \delta^{(2)}) \sigma'(z_1^{(1)}) = (0.7)(-0.5189)(0.2403) \approx -0.0873, \nonumber \\
\delta_2^{(1)} &= (w_2^{(2)} \delta^{(2)}) \sigma'(z_2^{(1)}) = (-0.3)(-0.5189)(0.2484) \approx +0.0387.
\end{align}

Step 3: Compute gradients for hidden weights and biases via Equations BP3 and BP4:
\begin{align}
\frac{\partial \mathcal{L}}{\partial \mathbf{W}^{(1)}} &= \boldsymbol{\delta}^{(1)} \mathbf{x}^\top = \begin{bmatrix} -0.0873 \\ +0.0387 \end{bmatrix} \begin{bmatrix} 0.5 & -0.2 \end{bmatrix} = \begin{bmatrix} -0.0437 & +0.0175 \\ +0.0194 & -0.0077 \end{bmatrix}, \nonumber \\
\frac{\partial \mathcal{L}}{\partial \mathbf{b}^{(1)}} &= \boldsymbol{\delta}^{(1)} = \begin{bmatrix} -0.0873 \\ +0.0387 \end{bmatrix}.
\end{align}
In a single backward sweep, all 9 parameter gradients are computed exactly at machine precision.

\paragraph{Computational and memory drawbacks of backpropagation and the mini-batch solution.}
Despite its foundational role, the backpropagation algorithm incurs two significant computational and physical bottlenecks:
\begin{enumerate}[noitemsep,topsep=2pt,parsep=0pt,partopsep=0pt]
  \item \textbf{The Computational Drawback of Full-Batch Evaluation}:
  In standard batch gradient descent, evaluating the true empirical risk gradient requires backpropagating through the entire training dataset:
  \begin{equation}
  \nabla_{\boldsymbol{\theta}} \hat{\mathcal{R}}_N(\boldsymbol{\theta}) = \frac{1}{N} \sum_{i=1}^N \nabla_{\boldsymbol{\theta}} \mathcal{L}_i(\boldsymbol{\theta}).
  \end{equation}
  For large datasets (e.g., ImageNet with $N = 1.3 \times 10^6$ images, or modern web-scale text corpora with $N \sim 10^9$--$10^{12}$ tokens), executing a single parameter update requires an intractable computational cost of $\mathcal{O}(N \cdot W)$ operations. Much of this computation is severely redundant, as gradients evaluated on similar examples within a massive dataset largely cancel or repeat directional information.

  \textbf{How Mini-Batching Resolves the Computational Drawback}:
  Mini-batch optimization resolves this computational barrier by approximating the full gradient using a small, randomly sampled subset $\mathcal{B} \subset \{1, \dots, N\}$ of size $B \ll N$ (typically $B \in \{32, 64, 128, 256, 512\}$):
  \begin{equation}
  \mathbf{g}_{\mathcal{B}}(\boldsymbol{\theta}) = \frac{1}{B} \sum_{i \in \mathcal{B}} \nabla_{\boldsymbol{\theta}} \mathcal{L}_i(\boldsymbol{\theta}).
  \end{equation}
  Because observations are sampled uniformly at random, the mini-batch gradient is an \textbf{unbiased estimator} of the true full gradient: $\mathbb{E}_{\mathcal{B}}[\mathbf{g}_{\mathcal{B}}(\boldsymbol{\theta})] = \nabla_{\boldsymbol{\theta}} \hat{\mathcal{R}}_N(\boldsymbol{\theta})$. This reduces the computational cost per update from $\mathcal{O}(N \cdot W)$ down to $\mathcal{O}(B \cdot W)$---an acceleration of several orders of magnitude. 
  
  Furthermore, mini-batching unlocks massive hardware efficiency: by stacking $B$ observations into contiguous memory buffers, vector operations are transformed into dense General Matrix Multiplications ($\operatorname{GEMM}$) that saturate the high-throughput Tensor Cores of modern GPUs and TPUs. Performing $N/B$ parameter updates per epoch drastically accelerates convergence relative to full-batch optimization.

  \item \textbf{The Memory Drawback: Activation Caching and GPU VRAM Blowup}:
  As revealed by Equation BP3 ($\frac{\partial \mathcal{L}}{\partial \mathbf{W}^{(l)}} = \boldsymbol{\delta}^{(l)} (\mathbf{a}^{(l-1)})^\top$), computing the weight gradient at layer $l$ strictly requires the forward activation matrix $\mathbf{A}^{(l-1)}$. Consequently, all intermediate activations computed during the forward pass \textbf{must be retained in high-bandwidth memory (GPU VRAM)} until the backward pass traverses back to layer $l$.
  
  This introduces a critical memory asymmetry between inference and training:
  \begin{itemize}[noitemsep,topsep=2pt,parsep=0pt,partopsep=0pt]
    \item \textbf{Inference Memory}: Intermediate activations can be discarded immediately once the next layer is computed. Peak memory is determined only by the largest single layer: $\mathcal{O}(B \cdot \max_l d_l)$.
    \item \textbf{Training Memory}: All intermediate activations across \textbf{all $L$ layers} must be stored concurrently:
    \begin{equation}
    M_{\text{activations}} = B \sum_{l=1}^L d_l \approx \mathcal{O}(L \cdot B \cdot d).
    \label{eq:activation_memory_scaling}
    \end{equation}
  \end{itemize}
  In ultra-deep networks (such as 152-layer ResNets or deep Transformers), activation memory vastly exceeds parameter memory, forming the primary physical bottleneck on training batch size and model depth. To mitigate this VRAM wall, modern systems implement \textbf{Activation Checkpointing}~\cite{chen2016training}, which discards activations during the forward pass and selectively recomputes them during the backward sweep, trading a $30\%$ increase in compute FLOPs for an $\mathcal{O}(\sqrt{L})$ reduction in memory footprint.
\end{enumerate}
